\documentclass[
  doc,
  floatsintext,
  longtable,
  a4paper,
  nolmodern,
  notxfonts,
  notimes,
  12pt,
  colorlinks=true,linkcolor=blue,citecolor=blue,urlcolor=blue]{apa7}

\usepackage{amsmath}
\usepackage{amssymb}

\geometry{inner=1in, outer=1in}
\fancyhfoffset[LE,RO]{0cm}


\usepackage[bidi=default]{babel}
\babelprovide[main,import]{spanish}


% get rid of language-specific shorthands (see #6817):
\let\LanguageShortHands\languageshorthands
\def\languageshorthands#1{}

\RequirePackage{longtable}
\RequirePackage{threeparttablex}

\makeatletter
\renewcommand{\paragraph}{\@startsection{paragraph}{4}{\parindent}%
	{0\baselineskip \@plus 0.2ex \@minus 0.2ex}%
	{-.5em}%
	{\normalfont\normalsize\bfseries\typesectitle}}

\renewcommand{\subparagraph}[1]{\@startsection{subparagraph}{5}{0.5em}%
	{0\baselineskip \@plus 0.2ex \@minus 0.2ex}%
	{-\z@\relax}%
	{\normalfont\normalsize\bfseries\itshape\hspace{\parindent}{#1}\textit{\addperi}}{\relax}}
\makeatother




\usepackage{longtable, booktabs, multirow, multicol, colortbl, hhline, caption, array, float, xpatch}
\usepackage{subcaption}


\renewcommand\thesubfigure{\Alph{subfigure}}
\setcounter{topnumber}{2}
\setcounter{bottomnumber}{2}
\setcounter{totalnumber}{4}
\renewcommand{\topfraction}{0.85}
\renewcommand{\bottomfraction}{0.85}
\renewcommand{\textfraction}{0.15}
\renewcommand{\floatpagefraction}{0.7}

\usepackage{tcolorbox}
\tcbuselibrary{listings,theorems, breakable, skins}
\usepackage{fontawesome5}

\definecolor{quarto-callout-color}{HTML}{909090}
\definecolor{quarto-callout-note-color}{HTML}{0758E5}
\definecolor{quarto-callout-important-color}{HTML}{CC1914}
\definecolor{quarto-callout-warning-color}{HTML}{EB9113}
\definecolor{quarto-callout-tip-color}{HTML}{00A047}
\definecolor{quarto-callout-caution-color}{HTML}{FC5300}
\definecolor{quarto-callout-color-frame}{HTML}{ACACAC}
\definecolor{quarto-callout-note-color-frame}{HTML}{4582EC}
\definecolor{quarto-callout-important-color-frame}{HTML}{D9534F}
\definecolor{quarto-callout-warning-color-frame}{HTML}{F0AD4E}
\definecolor{quarto-callout-tip-color-frame}{HTML}{02B875}
\definecolor{quarto-callout-caution-color-frame}{HTML}{FD7E14}

%\newlength\Oldarrayrulewidth
%\newlength\Oldtabcolsep


\usepackage{hyperref}



\usepackage{color}
\usepackage{fancyvrb}
\newcommand{\VerbBar}{|}
\newcommand{\VERB}{\Verb[commandchars=\\\{\}]}
\DefineVerbatimEnvironment{Highlighting}{Verbatim}{commandchars=\\\{\}}
% Add ',fontsize=\small' for more characters per line
\usepackage{framed}
\definecolor{shadecolor}{RGB}{241,243,245}
\newenvironment{Shaded}{\begin{snugshade}}{\end{snugshade}}
\newcommand{\AlertTok}[1]{\textcolor[rgb]{0.68,0.00,0.00}{#1}}
\newcommand{\AnnotationTok}[1]{\textcolor[rgb]{0.37,0.37,0.37}{#1}}
\newcommand{\AttributeTok}[1]{\textcolor[rgb]{0.40,0.45,0.13}{#1}}
\newcommand{\BaseNTok}[1]{\textcolor[rgb]{0.68,0.00,0.00}{#1}}
\newcommand{\BuiltInTok}[1]{\textcolor[rgb]{0.00,0.23,0.31}{#1}}
\newcommand{\CharTok}[1]{\textcolor[rgb]{0.13,0.47,0.30}{#1}}
\newcommand{\CommentTok}[1]{\textcolor[rgb]{0.37,0.37,0.37}{#1}}
\newcommand{\CommentVarTok}[1]{\textcolor[rgb]{0.37,0.37,0.37}{\textit{#1}}}
\newcommand{\ConstantTok}[1]{\textcolor[rgb]{0.56,0.35,0.01}{#1}}
\newcommand{\ControlFlowTok}[1]{\textcolor[rgb]{0.00,0.23,0.31}{\textbf{#1}}}
\newcommand{\DataTypeTok}[1]{\textcolor[rgb]{0.68,0.00,0.00}{#1}}
\newcommand{\DecValTok}[1]{\textcolor[rgb]{0.68,0.00,0.00}{#1}}
\newcommand{\DocumentationTok}[1]{\textcolor[rgb]{0.37,0.37,0.37}{\textit{#1}}}
\newcommand{\ErrorTok}[1]{\textcolor[rgb]{0.68,0.00,0.00}{#1}}
\newcommand{\ExtensionTok}[1]{\textcolor[rgb]{0.00,0.23,0.31}{#1}}
\newcommand{\FloatTok}[1]{\textcolor[rgb]{0.68,0.00,0.00}{#1}}
\newcommand{\FunctionTok}[1]{\textcolor[rgb]{0.28,0.35,0.67}{#1}}
\newcommand{\ImportTok}[1]{\textcolor[rgb]{0.00,0.46,0.62}{#1}}
\newcommand{\InformationTok}[1]{\textcolor[rgb]{0.37,0.37,0.37}{#1}}
\newcommand{\KeywordTok}[1]{\textcolor[rgb]{0.00,0.23,0.31}{\textbf{#1}}}
\newcommand{\NormalTok}[1]{\textcolor[rgb]{0.00,0.23,0.31}{#1}}
\newcommand{\OperatorTok}[1]{\textcolor[rgb]{0.37,0.37,0.37}{#1}}
\newcommand{\OtherTok}[1]{\textcolor[rgb]{0.00,0.23,0.31}{#1}}
\newcommand{\PreprocessorTok}[1]{\textcolor[rgb]{0.68,0.00,0.00}{#1}}
\newcommand{\RegionMarkerTok}[1]{\textcolor[rgb]{0.00,0.23,0.31}{#1}}
\newcommand{\SpecialCharTok}[1]{\textcolor[rgb]{0.37,0.37,0.37}{#1}}
\newcommand{\SpecialStringTok}[1]{\textcolor[rgb]{0.13,0.47,0.30}{#1}}
\newcommand{\StringTok}[1]{\textcolor[rgb]{0.13,0.47,0.30}{#1}}
\newcommand{\VariableTok}[1]{\textcolor[rgb]{0.07,0.07,0.07}{#1}}
\newcommand{\VerbatimStringTok}[1]{\textcolor[rgb]{0.13,0.47,0.30}{#1}}
\newcommand{\WarningTok}[1]{\textcolor[rgb]{0.37,0.37,0.37}{\textit{#1}}}

\providecommand{\tightlist}{%
  \setlength{\itemsep}{0pt}\setlength{\parskip}{0pt}}
\usepackage{longtable,booktabs,array}
\newcounter{none} % for unnumbered tables
\usepackage{calc} % for calculating minipage widths
% Correct order of tables after \paragraph or \subparagraph
\usepackage{etoolbox}
\makeatletter
\patchcmd\longtable{\par}{\if@noskipsec\mbox{}\fi\par}{}{}
\makeatother
% Allow footnotes in longtable head/foot
\IfFileExists{footnotehyper.sty}{\usepackage{footnotehyper}}{\usepackage{footnote}}
\makesavenoteenv{longtable}

\usepackage{graphicx}
\makeatletter
\newsavebox\pandoc@box
\newcommand*\pandocbounded[1]{% scales image to fit in text height/width
  \sbox\pandoc@box{#1}%
  \Gscale@div\@tempa{\textheight}{\dimexpr\ht\pandoc@box+\dp\pandoc@box\relax}%
  \Gscale@div\@tempb{\linewidth}{\wd\pandoc@box}%
  \ifdim\@tempb\p@<\@tempa\p@\let\@tempa\@tempb\fi% select the smaller of both
  \ifdim\@tempa\p@<\p@\scalebox{\@tempa}{\usebox\pandoc@box}%
  \else\usebox{\pandoc@box}%
  \fi%
}
% Set default figure placement to htbp
\def\fps@figure{htbp}
\makeatother







\usepackage{newtx}

\defaultfontfeatures{Scale=MatchLowercase}
\defaultfontfeatures[\rmfamily]{Ligatures=TeX,Scale=1}





\title{Objetos y tipos de datos en Python}


\shorttitle{OBJETOS PYTHON}


\usepackage{etoolbox}



\ccoppy{\textcopyright~2021}



\author{Edison Achalma}



\affiliation{
{Escuela Profesional de Economía, Universidad Nacional de San Cristóbal
de Huamanga}}




\leftheader{Achalma}

\date{2021-06-07}


\abstract{En Python, todo es un objeto y su comportamiento está
determinado por su tipo y mutabilidad. Este artículo explora los
conceptos fundamentales de objetos en Python, con énfasis en la
diferencia entre asignación por referencia y por valor, los principales
tipos de datos built-in y el manejo práctico de las colecciones más
utilizadas; listas (mutables y ordenadas), tuplas (inmutables y
eficientes en memoria) y diccionarios (mutables, no ordenados y basados
en claves-valor). Se presentan ejemplos claros sobre indexación, métodos
principales, evaluación de membresía, copia superficial y profunda, así
como consideraciones importantes de consumo de memoria y buenas
prácticas al trabajar con estas estructuras. }

\keywords{Objetos Python, Mutabilidad, Referencias vs copia}

\authornote{\par{\addORCIDlink{Edison Achalma}{0000-0001-6996-3364}} 
\par{ }
\par{   El autor no tiene conflictos de interés que revelar.    Los
roles de autor se clasificaron utilizando la taxonomía de roles de
colaborador (CRediT; https://credit.niso.org/) de la siguiente
manera:  Edison Achalma:   conceptualización, metodología, análisis
formal, investigación, recursos, curación de
datos, redacción, visualización, supervisión, administración del
proyecto}
\par{La correspondencia relativa a este artículo debe dirigirse a Edison
Achalma, Escuela Profesional de Economía, Universidad Nacional de San
Cristóbal de Huamanga, Portal Independencia N°
57, Ayacucho, AYA 5001, Perú, Email: \href{mailto:elmer.achalma.09@unsch.edu.pe}{elmer.achalma.09@unsch.edu.pe}}
}

\makeatletter
\let\endoldlt\endlongtable
\def\endlongtable{
\hline
\endoldlt
}
\makeatother

\urlstyle{same}



\hypersetup{
  pdftitle={Objetos y tipos de datos en Python},
  pdfauthor={Autor}
}
\makeatletter
\@ifpackageloaded{caption}{}{\usepackage{caption}}
\AtBeginDocument{%
\ifdefined\contentsname
  \renewcommand*\contentsname{Tabla de contenidos}
\else
  \newcommand\contentsname{Tabla de contenidos}
\fi
\ifdefined\listfigurename
  \renewcommand*\listfigurename{Lista de Figuras}
\else
  \newcommand\listfigurename{Lista de Figuras}
\fi
\ifdefined\listtablename
  \renewcommand*\listtablename{Lista de Tablas}
\else
  \newcommand\listtablename{Lista de Tablas}
\fi
\ifdefined\figurename
  \renewcommand*\figurename{Figura}
\else
  \newcommand\figurename{Figura}
\fi
\ifdefined\tablename
  \renewcommand*\tablename{Tabla}
\else
  \newcommand\tablename{Tabla}
\fi
}
\@ifpackageloaded{float}{}{\usepackage{float}}
\floatstyle{ruled}
\@ifundefined{c@chapter}{\newfloat{codelisting}{h}{lop}}{\newfloat{codelisting}{h}{lop}[chapter]}
\floatname{codelisting}{Listado}
\newcommand*\listoflistings{\listof{codelisting}{Lista de Listados}}
\makeatother
\makeatletter
\makeatother
\makeatletter
\@ifpackageloaded{caption}{}{\usepackage{caption}}
\@ifpackageloaded{subcaption}{}{\usepackage{subcaption}}
\makeatother
\makeatletter
\@ifpackageloaded{fontawesome5}{}{\usepackage{fontawesome5}}
\makeatother

% From https://tex.stackexchange.com/a/645996/211326
%%% apa7 doesn't want to add appendix section titles in the toc
%%% let's make it do it
\makeatletter
\xpatchcmd{\appendix}
  {\par}
  {\addcontentsline{toc}{section}{\@currentlabelname}\par}
  {}{}
\makeatother

%% Disable longtable counter
%% https://tex.stackexchange.com/a/248395/211326

\usepackage{etoolbox}

\makeatletter
\patchcmd{\LT@caption}
  {\bgroup}
  {\bgroup\global\LTpatch@captiontrue}
  {}{}
\patchcmd{\longtable}
  {\par}
  {\par\global\LTpatch@captionfalse}
  {}{}
\apptocmd{\endlongtable}
  {\ifLTpatch@caption\else\addtocounter{table}{-1}\fi}
  {}{}
\newif\ifLTpatch@caption
\makeatother

\begin{document}

\maketitle


\hypertarget{toc}{}
\tableofcontents
\newpage
\section[Introduction]{Objetos y tipos de datos en Python}

\setcounter{secnumdepth}{3}

\setlength\LTleft{0pt}




En esta tercera guía exploraremos las estructuras de datos fundamentales
de Python: listas, diccionarios y tuplas. Estas estructuras son
esenciales para organizar y manipular datos económicos de manera
eficiente.

\section{Listas}\label{listas}

Las listas son colecciones ordenadas y mutables de elementos en Python.
Son uno de los tipos de datos más versátiles y utilizados, especialmente
en análisis económico donde necesitamos trabajar con series de datos.

\subsection{Creación de listas}\label{creaciuxf3n-de-listas}

\begin{Shaded}
\begin{Highlighting}[]
\CommentTok{\# Lista de precios}
\NormalTok{precios }\OperatorTok{=}\NormalTok{ [}\FloatTok{10.50}\NormalTok{, }\FloatTok{12.30}\NormalTok{, }\FloatTok{15.80}\NormalTok{, }\FloatTok{9.90}\NormalTok{, }\FloatTok{11.20}\NormalTok{]}
\BuiltInTok{print}\NormalTok{(precios)}

\CommentTok{\# Lista de sectores económicos}
\NormalTok{sectores }\OperatorTok{=}\NormalTok{ [}\StringTok{\textquotesingle{}Agricultura\textquotesingle{}}\NormalTok{, }\StringTok{\textquotesingle{}Minería\textquotesingle{}}\NormalTok{, }\StringTok{\textquotesingle{}Manufactura\textquotesingle{}}\NormalTok{, }\StringTok{\textquotesingle{}Servicios\textquotesingle{}}\NormalTok{]}
\BuiltInTok{print}\NormalTok{(sectores)}

\CommentTok{\# Lista de años}
\NormalTok{anos }\OperatorTok{=}\NormalTok{ [}\DecValTok{2019}\NormalTok{, }\DecValTok{2020}\NormalTok{, }\DecValTok{2021}\NormalTok{, }\DecValTok{2022}\NormalTok{, }\DecValTok{2023}\NormalTok{]}
\BuiltInTok{print}\NormalTok{(anos)}

\CommentTok{\# Crear lista con la función list()}
\NormalTok{tasas\_interes }\OperatorTok{=} \BuiltInTok{list}\NormalTok{([}\FloatTok{0.05}\NormalTok{, }\FloatTok{0.06}\NormalTok{, }\FloatTok{0.055}\NormalTok{, }\FloatTok{0.065}\NormalTok{, }\FloatTok{0.07}\NormalTok{])}
\BuiltInTok{print}\NormalTok{(tasas\_interes)}

\CommentTok{\# Lista vacía}
\NormalTok{datos\_vacios }\OperatorTok{=}\NormalTok{ []}
\BuiltInTok{print}\NormalTok{(datos\_vacios)}
\end{Highlighting}
\end{Shaded}

\subsection{Listas heterogéneas}\label{listas-heteroguxe9neas}

Las listas pueden contener elementos de diferentes tipos:

\begin{Shaded}
\begin{Highlighting}[]
\CommentTok{\# Lista mixta con datos económicos}
\NormalTok{datos\_empresa }\OperatorTok{=}\NormalTok{ [}\StringTok{\textquotesingle{}TechCorp\textquotesingle{}}\NormalTok{, }\DecValTok{2015}\NormalTok{, }\FloatTok{1500000.50}\NormalTok{, }\VariableTok{True}\NormalTok{, [}\StringTok{\textquotesingle{}Lima\textquotesingle{}}\NormalTok{, }\StringTok{\textquotesingle{}Arequipa\textquotesingle{}}\NormalTok{]]}
\BuiltInTok{print}\NormalTok{(datos\_empresa)}
\CommentTok{\# [\textquotesingle{}TechCorp\textquotesingle{}, 2015, 1500000.50, True, [\textquotesingle{}Lima\textquotesingle{}, \textquotesingle{}Arequipa\textquotesingle{}]]}

\CommentTok{\# Estructura de datos compleja}
\NormalTok{indicador\_economico }\OperatorTok{=}\NormalTok{ [}
    \StringTok{\textquotesingle{}PIB\textquotesingle{}}\NormalTok{,}
    \DecValTok{242632}\NormalTok{,  }\CommentTok{\# Millones de USD}
    \DecValTok{2023}\NormalTok{,}
\NormalTok{    [}\FloatTok{2.4}\NormalTok{, }\FloatTok{4.0}\NormalTok{, }\OperatorTok{{-}}\FloatTok{11.0}\NormalTok{, }\FloatTok{13.3}\NormalTok{, }\FloatTok{2.7}\NormalTok{],  }\CommentTok{\# Tasas de crecimiento histórico}
\NormalTok{    \{}\StringTok{\textquotesingle{}moneda\textquotesingle{}}\NormalTok{: }\StringTok{\textquotesingle{}USD\textquotesingle{}}\NormalTok{, }\StringTok{\textquotesingle{}unidad\textquotesingle{}}\NormalTok{: }\StringTok{\textquotesingle{}millones\textquotesingle{}}\NormalTok{\}}
\NormalTok{]}
\end{Highlighting}
\end{Shaded}

\subsection{Indexación de listas}\label{indexaciuxf3n-de-listas}

Las listas en Python usan indexación de base cero (el primer elemento
está en la posición 0):

\begin{Shaded}
\begin{Highlighting}[]
\CommentTok{\# Datos de inflación mensual}
\NormalTok{inflacion\_mensual }\OperatorTok{=}\NormalTok{ [}\FloatTok{0.5}\NormalTok{, }\FloatTok{0.3}\NormalTok{, }\FloatTok{0.4}\NormalTok{, }\FloatTok{0.6}\NormalTok{, }\FloatTok{0.2}\NormalTok{, }\FloatTok{0.5}\NormalTok{, }\FloatTok{0.7}\NormalTok{, }\FloatTok{0.4}\NormalTok{, }\FloatTok{0.5}\NormalTok{, }\FloatTok{0.6}\NormalTok{, }\FloatTok{0.3}\NormalTok{, }\FloatTok{0.4}\NormalTok{]}

\CommentTok{\# Acceder a elementos individuales}
\NormalTok{enero }\OperatorTok{=}\NormalTok{ inflacion\_mensual[}\DecValTok{0}\NormalTok{]}
\BuiltInTok{print}\NormalTok{(}\SpecialStringTok{f"Inflación de enero: }\SpecialCharTok{\{}\NormalTok{enero}\SpecialCharTok{\}}\SpecialStringTok{\%"}\NormalTok{)  }\CommentTok{\# 0.5\%}

\NormalTok{febrero }\OperatorTok{=}\NormalTok{ inflacion\_mensual[}\DecValTok{1}\NormalTok{]}
\BuiltInTok{print}\NormalTok{(}\SpecialStringTok{f"Inflación de febrero: }\SpecialCharTok{\{}\NormalTok{febrero}\SpecialCharTok{\}}\SpecialStringTok{\%"}\NormalTok{)  }\CommentTok{\# 0.3\%}

\NormalTok{diciembre }\OperatorTok{=}\NormalTok{ inflacion\_mensual[}\DecValTok{11}\NormalTok{]}
\BuiltInTok{print}\NormalTok{(}\SpecialStringTok{f"Inflación de diciembre: }\SpecialCharTok{\{}\NormalTok{diciembre}\SpecialCharTok{\}}\SpecialStringTok{\%"}\NormalTok{)  }\CommentTok{\# 0.4\%}

\CommentTok{\# Indexación negativa (desde el final)}
\NormalTok{ultimo\_mes }\OperatorTok{=}\NormalTok{ inflacion\_mensual[}\OperatorTok{{-}}\DecValTok{1}\NormalTok{]}
\BuiltInTok{print}\NormalTok{(}\SpecialStringTok{f"Último mes: }\SpecialCharTok{\{}\NormalTok{ultimo\_mes}\SpecialCharTok{\}}\SpecialStringTok{\%"}\NormalTok{)  }\CommentTok{\# 0.4\%}

\NormalTok{penultimo\_mes }\OperatorTok{=}\NormalTok{ inflacion\_mensual[}\OperatorTok{{-}}\DecValTok{2}\NormalTok{]}
\BuiltInTok{print}\NormalTok{(}\SpecialStringTok{f"Penúltimo mes: }\SpecialCharTok{\{}\NormalTok{penultimo\_mes}\SpecialCharTok{\}}\SpecialStringTok{\%"}\NormalTok{)  }\CommentTok{\# 0.3\%}
\end{Highlighting}
\end{Shaded}

\subsection{Slicing de listas}\label{slicing-de-listas}

El slicing permite extraer sublistas:

\begin{Shaded}
\begin{Highlighting}[]
\CommentTok{\# PIB trimestral}
\NormalTok{pib\_trimestral }\OperatorTok{=}\NormalTok{ [}\DecValTok{50000}\NormalTok{, }\DecValTok{52000}\NormalTok{, }\DecValTok{51500}\NormalTok{, }\DecValTok{53000}\NormalTok{, }\DecValTok{54000}\NormalTok{, }\DecValTok{55500}\NormalTok{, }\DecValTok{56000}\NormalTok{, }\DecValTok{57200}\NormalTok{]}

\CommentTok{\# Sintaxis: lista[inicio:fin:paso]}
\CommentTok{\# El índice \textquotesingle{}fin\textquotesingle{} no se incluye}

\CommentTok{\# Primer trimestre del primer año}
\NormalTok{q1\_ano1 }\OperatorTok{=}\NormalTok{ pib\_trimestral[}\DecValTok{0}\NormalTok{]}
\BuiltInTok{print}\NormalTok{(}\SpecialStringTok{f"Q1 Año 1: }\SpecialCharTok{\{}\NormalTok{q1\_ano1}\SpecialCharTok{\}}\SpecialStringTok{"}\NormalTok{)}

\CommentTok{\# Todos los trimestres del primer año}
\NormalTok{ano1 }\OperatorTok{=}\NormalTok{ pib\_trimestral[}\DecValTok{0}\NormalTok{:}\DecValTok{4}\NormalTok{]}
\BuiltInTok{print}\NormalTok{(}\SpecialStringTok{f"Año 1: }\SpecialCharTok{\{}\NormalTok{ano1}\SpecialCharTok{\}}\SpecialStringTok{"}\NormalTok{)  }\CommentTok{\# [50000, 52000, 51500, 53000]}

\CommentTok{\# Todos los trimestres del segundo año}
\NormalTok{ano2 }\OperatorTok{=}\NormalTok{ pib\_trimestral[}\DecValTok{4}\NormalTok{:}\DecValTok{8}\NormalTok{]}
\BuiltInTok{print}\NormalTok{(}\SpecialStringTok{f"Año 2: }\SpecialCharTok{\{}\NormalTok{ano2}\SpecialCharTok{\}}\SpecialStringTok{"}\NormalTok{)  }\CommentTok{\# [54000, 55500, 56000, 57200]}

\CommentTok{\# Primeros tres trimestres}
\NormalTok{primeros\_tres }\OperatorTok{=}\NormalTok{ pib\_trimestral[:}\DecValTok{3}\NormalTok{]}
\BuiltInTok{print}\NormalTok{(}\SpecialStringTok{f"Primeros 3 trimestres: }\SpecialCharTok{\{}\NormalTok{primeros\_tres}\SpecialCharTok{\}}\SpecialStringTok{"}\NormalTok{)}

\CommentTok{\# Desde el cuarto trimestre hasta el final}
\NormalTok{desde\_cuarto }\OperatorTok{=}\NormalTok{ pib\_trimestral[}\DecValTok{3}\NormalTok{:]}
\BuiltInTok{print}\NormalTok{(}\SpecialStringTok{f"Desde Q4: }\SpecialCharTok{\{}\NormalTok{desde\_cuarto}\SpecialCharTok{\}}\SpecialStringTok{"}\NormalTok{)}

\CommentTok{\# Todos los elementos (copia de la lista)}
\NormalTok{todos }\OperatorTok{=}\NormalTok{ pib\_trimestral[:]}
\BuiltInTok{print}\NormalTok{(}\SpecialStringTok{f"Todos: }\SpecialCharTok{\{}\NormalTok{todos}\SpecialCharTok{\}}\SpecialStringTok{"}\NormalTok{)}

\CommentTok{\# Con paso: cada 2 elementos}
\NormalTok{trimestres\_pares }\OperatorTok{=}\NormalTok{ pib\_trimestral[::}\DecValTok{2}\NormalTok{]}
\BuiltInTok{print}\NormalTok{(}\SpecialStringTok{f"Trimestres impares (Q1, Q3): }\SpecialCharTok{\{}\NormalTok{trimestres\_pares}\SpecialCharTok{\}}\SpecialStringTok{"}\NormalTok{)}

\CommentTok{\# Invertir lista}
\NormalTok{invertido }\OperatorTok{=}\NormalTok{ pib\_trimestral[::}\OperatorTok{{-}}\DecValTok{1}\NormalTok{]}
\BuiltInTok{print}\NormalTok{(}\SpecialStringTok{f"Invertido: }\SpecialCharTok{\{}\NormalTok{invertido}\SpecialCharTok{\}}\SpecialStringTok{"}\NormalTok{)}
\end{Highlighting}
\end{Shaded}

\subsection{Modificación de
elementos}\label{modificaciuxf3n-de-elementos}

Las listas son mutables, podemos cambiar sus elementos:

\begin{Shaded}
\begin{Highlighting}[]
\CommentTok{\# Proyecciones de crecimiento}
\NormalTok{crecimiento }\OperatorTok{=}\NormalTok{ [}\FloatTok{2.5}\NormalTok{, }\FloatTok{3.0}\NormalTok{, }\FloatTok{3.2}\NormalTok{, }\FloatTok{2.8}\NormalTok{, }\FloatTok{3.5}\NormalTok{]}
\BuiltInTok{print}\NormalTok{(}\SpecialStringTok{f"Original: }\SpecialCharTok{\{}\NormalTok{crecimiento}\SpecialCharTok{\}}\SpecialStringTok{"}\NormalTok{)}

\CommentTok{\# Modificar un elemento}
\NormalTok{crecimiento[}\DecValTok{2}\NormalTok{] }\OperatorTok{=} \FloatTok{3.3}  \CommentTok{\# Actualizar proyección del tercer año}
\BuiltInTok{print}\NormalTok{(}\SpecialStringTok{f"Modificado: }\SpecialCharTok{\{}\NormalTok{crecimiento}\SpecialCharTok{\}}\SpecialStringTok{"}\NormalTok{)}

\CommentTok{\# Modificar múltiples elementos con slicing}
\NormalTok{crecimiento[}\DecValTok{3}\NormalTok{:}\DecValTok{5}\NormalTok{] }\OperatorTok{=}\NormalTok{ [}\FloatTok{2.9}\NormalTok{, }\FloatTok{3.6}\NormalTok{]  }\CommentTok{\# Actualizar últimos dos años}
\BuiltInTok{print}\NormalTok{(}\SpecialStringTok{f"Actualizado: }\SpecialCharTok{\{}\NormalTok{crecimiento}\SpecialCharTok{\}}\SpecialStringTok{"}\NormalTok{)}

\CommentTok{\# Modificar usando indexación negativa}
\NormalTok{crecimiento[}\OperatorTok{{-}}\DecValTok{1}\NormalTok{] }\OperatorTok{=} \FloatTok{3.7}  \CommentTok{\# Cambiar última proyección}
\BuiltInTok{print}\NormalTok{(}\SpecialStringTok{f"Final: }\SpecialCharTok{\{}\NormalTok{crecimiento}\SpecialCharTok{\}}\SpecialStringTok{"}\NormalTok{)}
\end{Highlighting}
\end{Shaded}

\subsection{Listas anidadas}\label{listas-anidadas}

Las listas pueden contener otras listas, útil para representar matrices
o datos multidimensionales:

\begin{Shaded}
\begin{Highlighting}[]
\CommentTok{\# Matriz de datos económicos por región y año}
\CommentTok{\# Filas: regiones, Columnas: años}
\NormalTok{pib\_regional }\OperatorTok{=}\NormalTok{ [}
\NormalTok{    [}\DecValTok{100000}\NormalTok{, }\DecValTok{105000}\NormalTok{, }\DecValTok{110000}\NormalTok{],  }\CommentTok{\# Lima}
\NormalTok{    [}\DecValTok{50000}\NormalTok{, }\DecValTok{52000}\NormalTok{, }\DecValTok{54500}\NormalTok{],     }\CommentTok{\# Arequipa}
\NormalTok{    [}\DecValTok{30000}\NormalTok{, }\DecValTok{31500}\NormalTok{, }\DecValTok{33000}\NormalTok{],     }\CommentTok{\# Cusco}
\NormalTok{    [}\DecValTok{45000}\NormalTok{, }\DecValTok{47000}\NormalTok{, }\DecValTok{49000}\NormalTok{]      }\CommentTok{\# La Libertad}
\NormalTok{]}

\CommentTok{\# Acceder a PIB de Lima en el segundo año}
\NormalTok{lima\_ano2 }\OperatorTok{=}\NormalTok{ pib\_regional[}\DecValTok{0}\NormalTok{][}\DecValTok{1}\NormalTok{]}
\BuiltInTok{print}\NormalTok{(}\SpecialStringTok{f"PIB Lima año 2: }\SpecialCharTok{\{}\NormalTok{lima\_ano2}\SpecialCharTok{\}}\SpecialStringTok{"}\NormalTok{)  }\CommentTok{\# 105000}

\CommentTok{\# Acceder a PIB de Cusco en el tercer año}
\NormalTok{cusco\_ano3 }\OperatorTok{=}\NormalTok{ pib\_regional[}\DecValTok{2}\NormalTok{][}\DecValTok{2}\NormalTok{]}
\BuiltInTok{print}\NormalTok{(}\SpecialStringTok{f"PIB Cusco año 3: }\SpecialCharTok{\{}\NormalTok{cusco\_ano3}\SpecialCharTok{\}}\SpecialStringTok{"}\NormalTok{)  }\CommentTok{\# 33000}

\CommentTok{\# Acceder a todos los datos de Arequipa}
\NormalTok{arequipa\_todos }\OperatorTok{=}\NormalTok{ pib\_regional[}\DecValTok{1}\NormalTok{]}
\BuiltInTok{print}\NormalTok{(}\SpecialStringTok{f"PIB Arequipa todos los años: }\SpecialCharTok{\{}\NormalTok{arequipa\_todos}\SpecialCharTok{\}}\SpecialStringTok{"}\NormalTok{)}

\CommentTok{\# Modificar un valor específico}
\NormalTok{pib\_regional[}\DecValTok{1}\NormalTok{][}\DecValTok{2}\NormalTok{] }\OperatorTok{=} \DecValTok{55000}  \CommentTok{\# Actualizar Arequipa año 3}
\BuiltInTok{print}\NormalTok{(}\SpecialStringTok{f"PIB Arequipa actualizado: }\SpecialCharTok{\{}\NormalTok{pib\_regional[}\DecValTok{1}\NormalTok{]}\SpecialCharTok{\}}\SpecialStringTok{"}\NormalTok{)}
\end{Highlighting}
\end{Shaded}

\section{Asignación por referencia vs por
valor}\label{asignaciuxf3n-por-referencia-vs-por-valor}

Este es uno de los conceptos más importantes y frecuentemente
malentendidos en Python. Entender la diferencia puede evitar errores
difíciles de detectar.

\subsection{Asignación por
referencia}\label{asignaciuxf3n-por-referencia}

Cuando asignamos una lista a otra variable, Python no crea una copia
nueva. Ambas variables apuntan al mismo objeto en memoria:

\begin{Shaded}
\begin{Highlighting}[]
\CommentTok{\# Lista original de precios}
\NormalTok{precios\_2023 }\OperatorTok{=}\NormalTok{ [}\DecValTok{100}\NormalTok{, }\DecValTok{105}\NormalTok{, }\DecValTok{110}\NormalTok{, }\DecValTok{108}\NormalTok{, }\DecValTok{112}\NormalTok{]}
\BuiltInTok{print}\NormalTok{(}\SpecialStringTok{f"Precios 2023: }\SpecialCharTok{\{}\NormalTok{precios\_2023}\SpecialCharTok{\}}\SpecialStringTok{"}\NormalTok{)}

\CommentTok{\# Asignación por referencia}
\NormalTok{precios\_2024 }\OperatorTok{=}\NormalTok{ precios\_2023  }\CommentTok{\# ¡No es una copia!}

\CommentTok{\# Modificar precios\_2024}
\NormalTok{precios\_2024[}\DecValTok{0}\NormalTok{] }\OperatorTok{=} \DecValTok{120}

\CommentTok{\# Ambas listas cambian}
\BuiltInTok{print}\NormalTok{(}\SpecialStringTok{f"Precios 2023: }\SpecialCharTok{\{}\NormalTok{precios\_2023}\SpecialCharTok{\}}\SpecialStringTok{"}\NormalTok{)  }\CommentTok{\# [120, 105, 110, 108, 112]}
\BuiltInTok{print}\NormalTok{(}\SpecialStringTok{f"Precios 2024: }\SpecialCharTok{\{}\NormalTok{precios\_2024}\SpecialCharTok{\}}\SpecialStringTok{"}\NormalTok{)  }\CommentTok{\# [120, 105, 110, 108, 112]}

\CommentTok{\# Verificar que son el mismo objeto}
\BuiltInTok{print}\NormalTok{(}\SpecialStringTok{f"¿Son el mismo objeto? }\SpecialCharTok{\{}\NormalTok{precios\_2023 }\KeywordTok{is}\NormalTok{ precios\_2024}\SpecialCharTok{\}}\SpecialStringTok{"}\NormalTok{)  }\CommentTok{\# True}
\end{Highlighting}
\end{Shaded}

Diagrama conceptual:

\begin{verbatim}
precios_2023 ──┐
               ├──> [120, 105, 110, 108, 112] (un solo objeto en memoria)
precios_2024 ──┘
\end{verbatim}

\subsection{Copia superficial (shallow
copy)}\label{copia-superficial-shallow-copy}

Para crear una copia independiente de una lista:

\begin{Shaded}
\begin{Highlighting}[]
\ImportTok{import}\NormalTok{ copy}

\CommentTok{\# Lista original}
\NormalTok{inflacion\_base }\OperatorTok{=}\NormalTok{ [}\FloatTok{2.5}\NormalTok{, }\FloatTok{3.0}\NormalTok{, }\FloatTok{2.8}\NormalTok{, }\FloatTok{3.2}\NormalTok{]}
\BuiltInTok{print}\NormalTok{(}\SpecialStringTok{f"Inflación base: }\SpecialCharTok{\{}\NormalTok{inflacion\_base}\SpecialCharTok{\}}\SpecialStringTok{"}\NormalTok{)}

\CommentTok{\# Método 1: usando copy.copy()}
\NormalTok{inflacion\_copia1 }\OperatorTok{=}\NormalTok{ copy.copy(inflacion\_base)}

\CommentTok{\# Método 2: usando slicing [:]}
\NormalTok{inflacion\_copia2 }\OperatorTok{=}\NormalTok{ inflacion\_base[:]}

\CommentTok{\# Método 3: usando list()}
\NormalTok{inflacion\_copia3 }\OperatorTok{=} \BuiltInTok{list}\NormalTok{(inflacion\_base)}

\CommentTok{\# Modificar la copia}
\NormalTok{inflacion\_copia1[}\DecValTok{0}\NormalTok{] }\OperatorTok{=} \FloatTok{2.7}

\CommentTok{\# Solo cambia la copia}
\BuiltInTok{print}\NormalTok{(}\SpecialStringTok{f"Inflación base:   }\SpecialCharTok{\{}\NormalTok{inflacion\_base}\SpecialCharTok{\}}\SpecialStringTok{"}\NormalTok{)   }\CommentTok{\# [2.5, 3.0, 2.8, 3.2]}
\BuiltInTok{print}\NormalTok{(}\SpecialStringTok{f"Inflación copia1: }\SpecialCharTok{\{}\NormalTok{inflacion\_copia1}\SpecialCharTok{\}}\SpecialStringTok{"}\NormalTok{) }\CommentTok{\# [2.7, 3.0, 2.8, 3.2]}

\CommentTok{\# Verificar que son objetos diferentes}
\BuiltInTok{print}\NormalTok{(}\SpecialStringTok{f"¿Son el mismo objeto? }\SpecialCharTok{\{}\NormalTok{inflacion\_base }\KeywordTok{is}\NormalTok{ inflacion\_copia1}\SpecialCharTok{\}}\SpecialStringTok{"}\NormalTok{)  }\CommentTok{\# False}
\end{Highlighting}
\end{Shaded}

\subsection{Problema con listas
anidadas}\label{problema-con-listas-anidadas}

La copia superficial no copia las listas internas:

\begin{Shaded}
\begin{Highlighting}[]
\ImportTok{import}\NormalTok{ copy}

\CommentTok{\# Lista de listas}
\NormalTok{datos\_economicos }\OperatorTok{=}\NormalTok{ [}
\NormalTok{    [}\StringTok{\textquotesingle{}Lima\textquotesingle{}}\NormalTok{, }\DecValTok{100000}\NormalTok{],}
\NormalTok{    [}\StringTok{\textquotesingle{}Arequipa\textquotesingle{}}\NormalTok{, }\DecValTok{50000}\NormalTok{],}
\NormalTok{    [}\StringTok{\textquotesingle{}Cusco\textquotesingle{}}\NormalTok{, }\DecValTok{30000}\NormalTok{]}
\NormalTok{]}

\CommentTok{\# Copia superficial}
\NormalTok{datos\_copia\_superficial }\OperatorTok{=}\NormalTok{ copy.copy(datos\_economicos)}

\CommentTok{\# Modificar una sublista}
\NormalTok{datos\_copia\_superficial[}\DecValTok{0}\NormalTok{][}\DecValTok{1}\NormalTok{] }\OperatorTok{=} \DecValTok{105000}

\CommentTok{\# La lista original también cambia}
\BuiltInTok{print}\NormalTok{(}\SpecialStringTok{f"Original: }\SpecialCharTok{\{}\NormalTok{datos\_economicos}\SpecialCharTok{\}}\SpecialStringTok{"}\NormalTok{)}
\CommentTok{\# [[\textquotesingle{}Lima\textquotesingle{}, 105000], [\textquotesingle{}Arequipa\textquotesingle{}, 50000], [\textquotesingle{}Cusco\textquotesingle{}, 30000]]}

\BuiltInTok{print}\NormalTok{(}\SpecialStringTok{f"Copia: }\SpecialCharTok{\{}\NormalTok{datos\_copia\_superficial}\SpecialCharTok{\}}\SpecialStringTok{"}\NormalTok{)}
\CommentTok{\# [[\textquotesingle{}Lima\textquotesingle{}, 105000], [\textquotesingle{}Arequipa\textquotesingle{}, 50000], [\textquotesingle{}Cusco\textquotesingle{}, 30000]]}
\end{Highlighting}
\end{Shaded}

\subsection{Copia profunda (deep copy)}\label{copia-profunda-deep-copy}

Para copiar listas que contienen otras listas:

\begin{Shaded}
\begin{Highlighting}[]
\ImportTok{import}\NormalTok{ copy}

\CommentTok{\# Lista de listas}
\NormalTok{datos\_economicos }\OperatorTok{=}\NormalTok{ [}
\NormalTok{    [}\StringTok{\textquotesingle{}Lima\textquotesingle{}}\NormalTok{, }\DecValTok{100000}\NormalTok{],}
\NormalTok{    [}\StringTok{\textquotesingle{}Arequipa\textquotesingle{}}\NormalTok{, }\DecValTok{50000}\NormalTok{],}
\NormalTok{    [}\StringTok{\textquotesingle{}Cusco\textquotesingle{}}\NormalTok{, }\DecValTok{30000}\NormalTok{]}
\NormalTok{]}

\CommentTok{\# Copia profunda}
\NormalTok{datos\_copia\_profunda }\OperatorTok{=}\NormalTok{ copy.deepcopy(datos\_economicos)}

\CommentTok{\# Modificar la copia}
\NormalTok{datos\_copia\_profunda[}\DecValTok{0}\NormalTok{][}\DecValTok{1}\NormalTok{] }\OperatorTok{=} \DecValTok{105000}

\CommentTok{\# Solo cambia la copia profunda}
\BuiltInTok{print}\NormalTok{(}\SpecialStringTok{f"Original: }\SpecialCharTok{\{}\NormalTok{datos\_economicos}\SpecialCharTok{\}}\SpecialStringTok{"}\NormalTok{)}
\CommentTok{\# [[\textquotesingle{}Lima\textquotesingle{}, 100000], [\textquotesingle{}Arequipa\textquotesingle{}, 50000], [\textquotesingle{}Cusco\textquotesingle{}, 30000]]}

\BuiltInTok{print}\NormalTok{(}\SpecialStringTok{f"Copia profunda: }\SpecialCharTok{\{}\NormalTok{datos\_copia\_profunda}\SpecialCharTok{\}}\SpecialStringTok{"}\NormalTok{)}
\CommentTok{\# [[\textquotesingle{}Lima\textquotesingle{}, 105000], [\textquotesingle{}Arequipa\textquotesingle{}, 50000], [\textquotesingle{}Cusco\textquotesingle{}, 30000]]}
\end{Highlighting}
\end{Shaded}

\subsection{Ejemplo aplicado: Escenarios
económicos}\label{ejemplo-aplicado-escenarios-econuxf3micos}

\begin{Shaded}
\begin{Highlighting}[]
\ImportTok{import}\NormalTok{ copy}

\CommentTok{\# Escenario base}
\NormalTok{escenario\_base }\OperatorTok{=}\NormalTok{ \{}
    \StringTok{\textquotesingle{}pib\textquotesingle{}}\NormalTok{: }\DecValTok{242632}\NormalTok{,}
    \StringTok{\textquotesingle{}inflacion\textquotesingle{}}\NormalTok{: }\FloatTok{3.5}\NormalTok{,}
    \StringTok{\textquotesingle{}desempleo\textquotesingle{}}\NormalTok{: }\FloatTok{7.2}\NormalTok{,}
    \StringTok{\textquotesingle{}sectores\textquotesingle{}}\NormalTok{: [}\StringTok{\textquotesingle{}Minería\textquotesingle{}}\NormalTok{, }\StringTok{\textquotesingle{}Agricultura\textquotesingle{}}\NormalTok{, }\StringTok{\textquotesingle{}Manufactura\textquotesingle{}}\NormalTok{]}
\NormalTok{\}}

\CommentTok{\# Crear escenarios alternativos}
\NormalTok{escenario\_optimista }\OperatorTok{=}\NormalTok{ copy.deepcopy(escenario\_base)}
\NormalTok{escenario\_pesimista }\OperatorTok{=}\NormalTok{ copy.deepcopy(escenario\_base)}

\CommentTok{\# Modificar escenarios}
\NormalTok{escenario\_optimista[}\StringTok{\textquotesingle{}pib\textquotesingle{}}\NormalTok{] }\OperatorTok{=} \DecValTok{250000}
\NormalTok{escenario\_optimista[}\StringTok{\textquotesingle{}inflacion\textquotesingle{}}\NormalTok{] }\OperatorTok{=} \FloatTok{3.0}
\NormalTok{escenario\_optimista[}\StringTok{\textquotesingle{}sectores\textquotesingle{}}\NormalTok{].append(}\StringTok{\textquotesingle{}Tecnología\textquotesingle{}}\NormalTok{)}

\NormalTok{escenario\_pesimista[}\StringTok{\textquotesingle{}pib\textquotesingle{}}\NormalTok{] }\OperatorTok{=} \DecValTok{235000}
\NormalTok{escenario\_pesimista[}\StringTok{\textquotesingle{}inflacion\textquotesingle{}}\NormalTok{] }\OperatorTok{=} \FloatTok{4.5}
\NormalTok{escenario\_pesimista[}\StringTok{\textquotesingle{}sectores\textquotesingle{}}\NormalTok{].remove(}\StringTok{\textquotesingle{}Manufactura\textquotesingle{}}\NormalTok{)}

\CommentTok{\# Cada escenario es independiente}
\BuiltInTok{print}\NormalTok{(}\StringTok{"Escenario Base:"}\NormalTok{)}
\BuiltInTok{print}\NormalTok{(escenario\_base)}
\BuiltInTok{print}\NormalTok{(}\StringTok{"}\CharTok{\textbackslash{}n}\StringTok{Escenario Optimista:"}\NormalTok{)}
\BuiltInTok{print}\NormalTok{(escenario\_optimista)}
\BuiltInTok{print}\NormalTok{(}\StringTok{"}\CharTok{\textbackslash{}n}\StringTok{Escenario Pesimista:"}\NormalTok{)}
\BuiltInTok{print}\NormalTok{(escenario\_pesimista)}
\end{Highlighting}
\end{Shaded}

\section{Operaciones con listas}\label{operaciones-con-listas}

\subsection{Concatenación y
repetición}\label{concatenaciuxf3n-y-repeticiuxf3n}

\begin{Shaded}
\begin{Highlighting}[]
\CommentTok{\# Datos del primer semestre}
\NormalTok{primer\_semestre }\OperatorTok{=}\NormalTok{ [}\DecValTok{50000}\NormalTok{, }\DecValTok{52000}\NormalTok{, }\DecValTok{51500}\NormalTok{]}

\CommentTok{\# Datos del segundo semestre}
\NormalTok{segundo\_semestre }\OperatorTok{=}\NormalTok{ [}\DecValTok{53000}\NormalTok{, }\DecValTok{54500}\NormalTok{, }\DecValTok{55000}\NormalTok{]}

\CommentTok{\# Concatenar listas}
\NormalTok{ano\_completo }\OperatorTok{=}\NormalTok{ primer\_semestre }\OperatorTok{+}\NormalTok{ segundo\_semestre}
\BuiltInTok{print}\NormalTok{(}\SpecialStringTok{f"Año completo: }\SpecialCharTok{\{}\NormalTok{ano\_completo}\SpecialCharTok{\}}\SpecialStringTok{"}\NormalTok{)}
\CommentTok{\# [50000, 52000, 51500, 53000, 54500, 55000]}

\CommentTok{\# Repetir listas}
\NormalTok{patron\_base }\OperatorTok{=}\NormalTok{ [}\DecValTok{100}\NormalTok{, }\DecValTok{110}\NormalTok{]}
\NormalTok{patron\_repetido }\OperatorTok{=}\NormalTok{ patron\_base }\OperatorTok{*} \DecValTok{3}
\BuiltInTok{print}\NormalTok{(}\SpecialStringTok{f"Patrón repetido: }\SpecialCharTok{\{}\NormalTok{patron\_repetido}\SpecialCharTok{\}}\SpecialStringTok{"}\NormalTok{)}
\CommentTok{\# [100, 110, 100, 110, 100, 110]}

\CommentTok{\# Aplicación: Crear valores iniciales repetidos}
\NormalTok{tasas\_iniciales }\OperatorTok{=}\NormalTok{ [}\FloatTok{0.05}\NormalTok{] }\OperatorTok{*} \DecValTok{12}  \CommentTok{\# Tasa fija para 12 meses}
\BuiltInTok{print}\NormalTok{(}\SpecialStringTok{f"Tasas iniciales: }\SpecialCharTok{\{}\NormalTok{tasas\_iniciales}\SpecialCharTok{\}}\SpecialStringTok{"}\NormalTok{)}
\end{Highlighting}
\end{Shaded}

\subsection{Pertenencia a listas}\label{pertenencia-a-listas}

\begin{Shaded}
\begin{Highlighting}[]
\CommentTok{\# Sectores económicos}
\NormalTok{sectores\_primarios }\OperatorTok{=}\NormalTok{ [}\StringTok{\textquotesingle{}Agricultura\textquotesingle{}}\NormalTok{, }\StringTok{\textquotesingle{}Pesca\textquotesingle{}}\NormalTok{, }\StringTok{\textquotesingle{}Minería\textquotesingle{}}\NormalTok{]}

\CommentTok{\# Verificar pertenencia}
\BuiltInTok{print}\NormalTok{(}\StringTok{\textquotesingle{}Minería\textquotesingle{}} \KeywordTok{in}\NormalTok{ sectores\_primarios)        }\CommentTok{\# True}
\BuiltInTok{print}\NormalTok{(}\StringTok{\textquotesingle{}Manufactura\textquotesingle{}} \KeywordTok{in}\NormalTok{ sectores\_primarios)    }\CommentTok{\# False}
\BuiltInTok{print}\NormalTok{(}\StringTok{\textquotesingle{}Pesca\textquotesingle{}} \KeywordTok{not} \KeywordTok{in}\NormalTok{ sectores\_primarios)      }\CommentTok{\# False}
\BuiltInTok{print}\NormalTok{(}\StringTok{\textquotesingle{}Tecnología\textquotesingle{}} \KeywordTok{not} \KeywordTok{in}\NormalTok{ sectores\_primarios) }\CommentTok{\# True}

\CommentTok{\# Aplicación práctica}
\NormalTok{sector\_analizar }\OperatorTok{=} \StringTok{\textquotesingle{}Agricultura\textquotesingle{}}

\ControlFlowTok{if}\NormalTok{ sector\_analizar }\KeywordTok{in}\NormalTok{ sectores\_primarios:}
    \BuiltInTok{print}\NormalTok{(}\SpecialStringTok{f"}\SpecialCharTok{\{}\NormalTok{sector\_analizar}\SpecialCharTok{\}}\SpecialStringTok{ es un sector primario"}\NormalTok{)}
\ControlFlowTok{else}\NormalTok{:}
    \BuiltInTok{print}\NormalTok{(}\SpecialStringTok{f"}\SpecialCharTok{\{}\NormalTok{sector\_analizar}\SpecialCharTok{\}}\SpecialStringTok{ no es un sector primario"}\NormalTok{)}
\end{Highlighting}
\end{Shaded}

\subsection{Comparación de listas}\label{comparaciuxf3n-de-listas}

\begin{Shaded}
\begin{Highlighting}[]
\CommentTok{\# PIB por trimestre de dos años}
\NormalTok{pib\_2022 }\OperatorTok{=}\NormalTok{ [}\DecValTok{50000}\NormalTok{, }\DecValTok{52000}\NormalTok{, }\DecValTok{51500}\NormalTok{, }\DecValTok{53000}\NormalTok{]}
\NormalTok{pib\_2023 }\OperatorTok{=}\NormalTok{ [}\DecValTok{54000}\NormalTok{, }\DecValTok{55000}\NormalTok{, }\DecValTok{54500}\NormalTok{, }\DecValTok{56000}\NormalTok{]}
\NormalTok{pib\_2023\_copia }\OperatorTok{=}\NormalTok{ [}\DecValTok{54000}\NormalTok{, }\DecValTok{55000}\NormalTok{, }\DecValTok{54500}\NormalTok{, }\DecValTok{56000}\NormalTok{]}

\CommentTok{\# Igualdad}
\BuiltInTok{print}\NormalTok{(pib\_2022 }\OperatorTok{==}\NormalTok{ pib\_2023)        }\CommentTok{\# False}
\BuiltInTok{print}\NormalTok{(pib\_2023 }\OperatorTok{==}\NormalTok{ pib\_2023\_copia)  }\CommentTok{\# True}

\CommentTok{\# Identidad (mismo objeto en memoria)}
\BuiltInTok{print}\NormalTok{(pib\_2023 }\KeywordTok{is}\NormalTok{ pib\_2023\_copia)  }\CommentTok{\# False}

\CommentTok{\# Comparación elemento por elemento}
\BuiltInTok{print}\NormalTok{(pib\_2022 }\OperatorTok{\textless{}}\NormalTok{ pib\_2023)  }\CommentTok{\# True (compara primer elemento diferente)}
\end{Highlighting}
\end{Shaded}

\subsection{Eliminar elementos}\label{eliminar-elementos}

\begin{Shaded}
\begin{Highlighting}[]
\CommentTok{\# Lista de sectores}
\NormalTok{sectores }\OperatorTok{=}\NormalTok{ [}\StringTok{\textquotesingle{}Agricultura\textquotesingle{}}\NormalTok{, }\StringTok{\textquotesingle{}Minería\textquotesingle{}}\NormalTok{, }\StringTok{\textquotesingle{}Manufactura\textquotesingle{}}\NormalTok{, }\StringTok{\textquotesingle{}Construcción\textquotesingle{}}\NormalTok{, }\StringTok{\textquotesingle{}Servicios\textquotesingle{}}\NormalTok{]}
\BuiltInTok{print}\NormalTok{(}\SpecialStringTok{f"Original: }\SpecialCharTok{\{}\NormalTok{sectores}\SpecialCharTok{\}}\SpecialStringTok{"}\NormalTok{)}

\CommentTok{\# Método 1: del (eliminar por índice)}
\KeywordTok{del}\NormalTok{ sectores[}\DecValTok{2}\NormalTok{]  }\CommentTok{\# Elimina \textquotesingle{}Manufactura\textquotesingle{}}
\BuiltInTok{print}\NormalTok{(}\SpecialStringTok{f"Después de del: }\SpecialCharTok{\{}\NormalTok{sectores}\SpecialCharTok{\}}\SpecialStringTok{"}\NormalTok{)}

\CommentTok{\# Método 2: remove() (eliminar por valor)}
\NormalTok{sectores.remove(}\StringTok{\textquotesingle{}Construcción\textquotesingle{}}\NormalTok{)}
\BuiltInTok{print}\NormalTok{(}\SpecialStringTok{f"Después de remove: }\SpecialCharTok{\{}\NormalTok{sectores}\SpecialCharTok{\}}\SpecialStringTok{"}\NormalTok{)}

\CommentTok{\# Método 3: pop() (eliminar y retornar elemento)}
\NormalTok{sector\_eliminado }\OperatorTok{=}\NormalTok{ sectores.pop(}\DecValTok{1}\NormalTok{)  }\CommentTok{\# Elimina y retorna \textquotesingle{}Minería\textquotesingle{}}
\BuiltInTok{print}\NormalTok{(}\SpecialStringTok{f"Sector eliminado: }\SpecialCharTok{\{}\NormalTok{sector\_eliminado}\SpecialCharTok{\}}\SpecialStringTok{"}\NormalTok{)}
\BuiltInTok{print}\NormalTok{(}\SpecialStringTok{f"Lista actual: }\SpecialCharTok{\{}\NormalTok{sectores}\SpecialCharTok{\}}\SpecialStringTok{"}\NormalTok{)}

\CommentTok{\# pop() sin argumentos elimina el último elemento}
\NormalTok{ultimo }\OperatorTok{=}\NormalTok{ sectores.pop()}
\BuiltInTok{print}\NormalTok{(}\SpecialStringTok{f"Último eliminado: }\SpecialCharTok{\{}\NormalTok{ultimo}\SpecialCharTok{\}}\SpecialStringTok{"}\NormalTok{)}
\BuiltInTok{print}\NormalTok{(}\SpecialStringTok{f"Lista final: }\SpecialCharTok{\{}\NormalTok{sectores}\SpecialCharTok{\}}\SpecialStringTok{"}\NormalTok{)}

\CommentTok{\# Método 4: clear() (vaciar toda la lista)}
\NormalTok{sectores.clear()}
\BuiltInTok{print}\NormalTok{(}\SpecialStringTok{f"Lista vacía: }\SpecialCharTok{\{}\NormalTok{sectores}\SpecialCharTok{\}}\SpecialStringTok{"}\NormalTok{)}
\end{Highlighting}
\end{Shaded}

\section{Métodos de listas}\label{muxe9todos-de-listas}

\subsection{Funciones básicas}\label{funciones-buxe1sicas}

\begin{Shaded}
\begin{Highlighting}[]
\CommentTok{\# Datos de ventas mensuales (en miles)}
\NormalTok{ventas }\OperatorTok{=}\NormalTok{ [}\DecValTok{45}\NormalTok{, }\DecValTok{52}\NormalTok{, }\DecValTok{48}\NormalTok{, }\DecValTok{55}\NormalTok{, }\DecValTok{60}\NormalTok{, }\DecValTok{58}\NormalTok{, }\DecValTok{62}\NormalTok{, }\DecValTok{59}\NormalTok{, }\DecValTok{65}\NormalTok{, }\DecValTok{70}\NormalTok{, }\DecValTok{68}\NormalTok{, }\DecValTok{75}\NormalTok{]}

\CommentTok{\# Longitud de la lista}
\NormalTok{num\_meses }\OperatorTok{=} \BuiltInTok{len}\NormalTok{(ventas)}
\BuiltInTok{print}\NormalTok{(}\SpecialStringTok{f"Número de meses: }\SpecialCharTok{\{}\NormalTok{num\_meses}\SpecialCharTok{\}}\SpecialStringTok{"}\NormalTok{)  }\CommentTok{\# 12}

\CommentTok{\# Valor mínimo}
\NormalTok{venta\_minima }\OperatorTok{=} \BuiltInTok{min}\NormalTok{(ventas)}
\BuiltInTok{print}\NormalTok{(}\SpecialStringTok{f"Venta mínima: }\SpecialCharTok{\{}\NormalTok{venta\_minima}\SpecialCharTok{\}}\SpecialStringTok{ mil"}\NormalTok{)  }\CommentTok{\# 45}

\CommentTok{\# Valor máximo}
\NormalTok{venta\_maxima }\OperatorTok{=} \BuiltInTok{max}\NormalTok{(ventas)}
\BuiltInTok{print}\NormalTok{(}\SpecialStringTok{f"Venta máxima: }\SpecialCharTok{\{}\NormalTok{venta\_maxima}\SpecialCharTok{\}}\SpecialStringTok{ mil"}\NormalTok{)  }\CommentTok{\# 75}

\CommentTok{\# Suma total}
\NormalTok{ventas\_totales }\OperatorTok{=} \BuiltInTok{sum}\NormalTok{(ventas)}
\BuiltInTok{print}\NormalTok{(}\SpecialStringTok{f"Ventas totales: }\SpecialCharTok{\{}\NormalTok{ventas\_totales}\SpecialCharTok{\}}\SpecialStringTok{ mil"}\NormalTok{)  }\CommentTok{\# 717}

\CommentTok{\# Promedio}
\NormalTok{ventas\_promedio }\OperatorTok{=} \BuiltInTok{sum}\NormalTok{(ventas) }\OperatorTok{/} \BuiltInTok{len}\NormalTok{(ventas)}
\BuiltInTok{print}\NormalTok{(}\SpecialStringTok{f"Venta promedio: }\SpecialCharTok{\{}\NormalTok{ventas\_promedio}\SpecialCharTok{:.2f\}}\SpecialStringTok{ mil"}\NormalTok{)  }\CommentTok{\# 59.75}
\end{Highlighting}
\end{Shaded}

\subsection{Método append()}\label{muxe9todo-append}

Agrega un elemento al final de la lista:

\begin{Shaded}
\begin{Highlighting}[]
\CommentTok{\# Registro de inflación mensual}
\NormalTok{inflacion }\OperatorTok{=}\NormalTok{ [}\FloatTok{0.5}\NormalTok{, }\FloatTok{0.3}\NormalTok{, }\FloatTok{0.4}\NormalTok{]}
\BuiltInTok{print}\NormalTok{(}\SpecialStringTok{f"Inflación inicial: }\SpecialCharTok{\{}\NormalTok{inflacion}\SpecialCharTok{\}}\SpecialStringTok{"}\NormalTok{)}

\CommentTok{\# Agregar nuevos meses}
\NormalTok{inflacion.append(}\FloatTok{0.6}\NormalTok{)}
\NormalTok{inflacion.append(}\FloatTok{0.5}\NormalTok{)}
\BuiltInTok{print}\NormalTok{(}\SpecialStringTok{f"Inflación actualizada: }\SpecialCharTok{\{}\NormalTok{inflacion}\SpecialCharTok{\}}\SpecialStringTok{"}\NormalTok{)}

\CommentTok{\# Agregar una lista como elemento único}
\NormalTok{inflacion.append([}\FloatTok{0.7}\NormalTok{, }\FloatTok{0.8}\NormalTok{])}
\BuiltInTok{print}\NormalTok{(}\SpecialStringTok{f"Con lista anidada: }\SpecialCharTok{\{}\NormalTok{inflacion}\SpecialCharTok{\}}\SpecialStringTok{"}\NormalTok{)}
\CommentTok{\# [0.5, 0.3, 0.4, 0.6, 0.5, [0.7, 0.8]]}
\end{Highlighting}
\end{Shaded}

\subsection{Método extend()}\label{muxe9todo-extend}

Extiende la lista agregando elementos de un iterable:

\begin{Shaded}
\begin{Highlighting}[]
\CommentTok{\# PIB trimestral}
\NormalTok{pib\_q1\_q2 }\OperatorTok{=}\NormalTok{ [}\DecValTok{50000}\NormalTok{, }\DecValTok{52000}\NormalTok{]}
\BuiltInTok{print}\NormalTok{(}\SpecialStringTok{f"PIB Q1{-}Q2: }\SpecialCharTok{\{}\NormalTok{pib\_q1\_q2}\SpecialCharTok{\}}\SpecialStringTok{"}\NormalTok{)}

\CommentTok{\# Extender con más trimestres}
\NormalTok{pib\_q3\_q4 }\OperatorTok{=}\NormalTok{ [}\DecValTok{51500}\NormalTok{, }\DecValTok{53000}\NormalTok{]}
\NormalTok{pib\_q1\_q2.extend(pib\_q3\_q4)}
\BuiltInTok{print}\NormalTok{(}\SpecialStringTok{f"PIB año completo: }\SpecialCharTok{\{}\NormalTok{pib\_q1\_q2}\SpecialCharTok{\}}\SpecialStringTok{"}\NormalTok{)}
\CommentTok{\# [50000, 52000, 51500, 53000]}

\CommentTok{\# Diferencia entre extend y append}
\NormalTok{lista1 }\OperatorTok{=}\NormalTok{ [}\DecValTok{1}\NormalTok{, }\DecValTok{2}\NormalTok{, }\DecValTok{3}\NormalTok{]}
\NormalTok{lista2 }\OperatorTok{=}\NormalTok{ [}\DecValTok{1}\NormalTok{, }\DecValTok{2}\NormalTok{, }\DecValTok{3}\NormalTok{]}

\NormalTok{lista1.extend([}\DecValTok{4}\NormalTok{, }\DecValTok{5}\NormalTok{])}
\NormalTok{lista2.append([}\DecValTok{4}\NormalTok{, }\DecValTok{5}\NormalTok{])}

\BuiltInTok{print}\NormalTok{(}\SpecialStringTok{f"Con extend: }\SpecialCharTok{\{}\NormalTok{lista1}\SpecialCharTok{\}}\SpecialStringTok{"}\NormalTok{)  }\CommentTok{\# [1, 2, 3, 4, 5]}
\BuiltInTok{print}\NormalTok{(}\SpecialStringTok{f"Con append: }\SpecialCharTok{\{}\NormalTok{lista2}\SpecialCharTok{\}}\SpecialStringTok{"}\NormalTok{)  }\CommentTok{\# [1, 2, 3, [4, 5]]}
\end{Highlighting}
\end{Shaded}

\subsection{Método insert()}\label{muxe9todo-insert}

Inserta un elemento en una posición específica:

\begin{Shaded}
\begin{Highlighting}[]
\CommentTok{\# Tasas de crecimiento}
\NormalTok{crecimiento }\OperatorTok{=}\NormalTok{ [}\FloatTok{2.5}\NormalTok{, }\FloatTok{3.0}\NormalTok{, }\FloatTok{3.5}\NormalTok{, }\FloatTok{4.0}\NormalTok{]}
\BuiltInTok{print}\NormalTok{(}\SpecialStringTok{f"Original: }\SpecialCharTok{\{}\NormalTok{crecimiento}\SpecialCharTok{\}}\SpecialStringTok{"}\NormalTok{)}

\CommentTok{\# Insertar en posición específica}
\NormalTok{crecimiento.insert(}\DecValTok{2}\NormalTok{, }\FloatTok{3.2}\NormalTok{)  }\CommentTok{\# Insertar 3.2 en índice 2}
\BuiltInTok{print}\NormalTok{(}\SpecialStringTok{f"Después de insert: }\SpecialCharTok{\{}\NormalTok{crecimiento}\SpecialCharTok{\}}\SpecialStringTok{"}\NormalTok{)}
\CommentTok{\# [2.5, 3.0, 3.2, 3.5, 4.0]}

\CommentTok{\# Insertar al inicio}
\NormalTok{crecimiento.insert(}\DecValTok{0}\NormalTok{, }\FloatTok{2.0}\NormalTok{)}
\BuiltInTok{print}\NormalTok{(}\SpecialStringTok{f"Con valor inicial: }\SpecialCharTok{\{}\NormalTok{crecimiento}\SpecialCharTok{\}}\SpecialStringTok{"}\NormalTok{)}
\CommentTok{\# [2.0, 2.5, 3.0, 3.2, 3.5, 4.0]}
\end{Highlighting}
\end{Shaded}

\subsection{Método remove()}\label{muxe9todo-remove}

Elimina la primera ocurrencia de un valor:

\begin{Shaded}
\begin{Highlighting}[]
\CommentTok{\# Sectores para análisis}
\NormalTok{sectores }\OperatorTok{=}\NormalTok{ [}\StringTok{\textquotesingle{}Minería\textquotesingle{}}\NormalTok{, }\StringTok{\textquotesingle{}Agricultura\textquotesingle{}}\NormalTok{, }\StringTok{\textquotesingle{}Minería\textquotesingle{}}\NormalTok{, }\StringTok{\textquotesingle{}Manufactura\textquotesingle{}}\NormalTok{]}
\BuiltInTok{print}\NormalTok{(}\SpecialStringTok{f"Original: }\SpecialCharTok{\{}\NormalTok{sectores}\SpecialCharTok{\}}\SpecialStringTok{"}\NormalTok{)}

\CommentTok{\# Eliminar primera ocurrencia de \textquotesingle{}Minería\textquotesingle{}}
\NormalTok{sectores.remove(}\StringTok{\textquotesingle{}Minería\textquotesingle{}}\NormalTok{)}
\BuiltInTok{print}\NormalTok{(}\SpecialStringTok{f"Después de remove: }\SpecialCharTok{\{}\NormalTok{sectores}\SpecialCharTok{\}}\SpecialStringTok{"}\NormalTok{)}
\CommentTok{\# [\textquotesingle{}Agricultura\textquotesingle{}, \textquotesingle{}Minería\textquotesingle{}, \textquotesingle{}Manufactura\textquotesingle{}]}

\CommentTok{\# Intento de eliminar valor inexistente genera error}
\ControlFlowTok{try}\NormalTok{:}
\NormalTok{    sectores.remove(}\StringTok{\textquotesingle{}Tecnología\textquotesingle{}}\NormalTok{)}
\ControlFlowTok{except} \PreprocessorTok{ValueError}\NormalTok{:}
    \BuiltInTok{print}\NormalTok{(}\StringTok{"Error: \textquotesingle{}Tecnología\textquotesingle{} no está en la lista"}\NormalTok{)}
\end{Highlighting}
\end{Shaded}

\subsection{Método pop()}\label{muxe9todo-pop}

Elimina y retorna un elemento:

\begin{Shaded}
\begin{Highlighting}[]
\CommentTok{\# Cola de proyectos}
\NormalTok{proyectos }\OperatorTok{=}\NormalTok{ [}\StringTok{\textquotesingle{}Proyecto A\textquotesingle{}}\NormalTok{, }\StringTok{\textquotesingle{}Proyecto B\textquotesingle{}}\NormalTok{, }\StringTok{\textquotesingle{}Proyecto C\textquotesingle{}}\NormalTok{, }\StringTok{\textquotesingle{}Proyecto D\textquotesingle{}}\NormalTok{]}
\BuiltInTok{print}\NormalTok{(}\SpecialStringTok{f"Proyectos: }\SpecialCharTok{\{}\NormalTok{proyectos}\SpecialCharTok{\}}\SpecialStringTok{"}\NormalTok{)}

\CommentTok{\# Eliminar y obtener el último proyecto}
\NormalTok{ultimo\_proyecto }\OperatorTok{=}\NormalTok{ proyectos.pop()}
\BuiltInTok{print}\NormalTok{(}\SpecialStringTok{f"Proyecto completado: }\SpecialCharTok{\{}\NormalTok{ultimo\_proyecto}\SpecialCharTok{\}}\SpecialStringTok{"}\NormalTok{)}
\BuiltInTok{print}\NormalTok{(}\SpecialStringTok{f"Proyectos restantes: }\SpecialCharTok{\{}\NormalTok{proyectos}\SpecialCharTok{\}}\SpecialStringTok{"}\NormalTok{)}

\CommentTok{\# Eliminar y obtener un proyecto específico}
\NormalTok{proyecto\_siguiente }\OperatorTok{=}\NormalTok{ proyectos.pop(}\DecValTok{0}\NormalTok{)  }\CommentTok{\# Primer proyecto}
\BuiltInTok{print}\NormalTok{(}\SpecialStringTok{f"Proyecto en ejecución: }\SpecialCharTok{\{}\NormalTok{proyecto\_siguiente}\SpecialCharTok{\}}\SpecialStringTok{"}\NormalTok{)}
\BuiltInTok{print}\NormalTok{(}\SpecialStringTok{f"Proyectos pendientes: }\SpecialCharTok{\{}\NormalTok{proyectos}\SpecialCharTok{\}}\SpecialStringTok{"}\NormalTok{)}
\end{Highlighting}
\end{Shaded}

\subsection{Método index()}\label{muxe9todo-index}

Retorna el índice de la primera ocurrencia de un valor:

\begin{Shaded}
\begin{Highlighting}[]
\CommentTok{\# Sectores económicos}
\NormalTok{sectores }\OperatorTok{=}\NormalTok{ [}\StringTok{\textquotesingle{}Agricultura\textquotesingle{}}\NormalTok{, }\StringTok{\textquotesingle{}Minería\textquotesingle{}}\NormalTok{, }\StringTok{\textquotesingle{}Manufactura\textquotesingle{}}\NormalTok{, }\StringTok{\textquotesingle{}Servicios\textquotesingle{}}\NormalTok{, }\StringTok{\textquotesingle{}Construcción\textquotesingle{}}\NormalTok{]}

\CommentTok{\# Encontrar índice de un sector}
\NormalTok{indice\_mineria }\OperatorTok{=}\NormalTok{ sectores.index(}\StringTok{\textquotesingle{}Minería\textquotesingle{}}\NormalTok{)}
\BuiltInTok{print}\NormalTok{(}\SpecialStringTok{f"Índice de Minería: }\SpecialCharTok{\{}\NormalTok{indice\_mineria}\SpecialCharTok{\}}\SpecialStringTok{"}\NormalTok{)  }\CommentTok{\# 1}

\NormalTok{indice\_servicios }\OperatorTok{=}\NormalTok{ sectores.index(}\StringTok{\textquotesingle{}Servicios\textquotesingle{}}\NormalTok{)}
\BuiltInTok{print}\NormalTok{(}\SpecialStringTok{f"Índice de Servicios: }\SpecialCharTok{\{}\NormalTok{indice\_servicios}\SpecialCharTok{\}}\SpecialStringTok{"}\NormalTok{)  }\CommentTok{\# 3}

\CommentTok{\# Búsqueda en rango específico}
\NormalTok{lista\_numeros }\OperatorTok{=}\NormalTok{ [}\DecValTok{1}\NormalTok{, }\DecValTok{2}\NormalTok{, }\DecValTok{3}\NormalTok{, }\DecValTok{2}\NormalTok{, }\DecValTok{4}\NormalTok{, }\DecValTok{2}\NormalTok{, }\DecValTok{5}\NormalTok{]}
\NormalTok{indice }\OperatorTok{=}\NormalTok{ lista\_numeros.index(}\DecValTok{2}\NormalTok{, }\DecValTok{2}\NormalTok{)  }\CommentTok{\# Buscar 2 desde índice 2}
\BuiltInTok{print}\NormalTok{(}\SpecialStringTok{f"Índice de 2 desde posición 2: }\SpecialCharTok{\{}\NormalTok{indice}\SpecialCharTok{\}}\SpecialStringTok{"}\NormalTok{)  }\CommentTok{\# 3}
\end{Highlighting}
\end{Shaded}

\subsection{Método count()}\label{muxe9todo-count}

Cuenta las ocurrencias de un valor:

\begin{Shaded}
\begin{Highlighting}[]
\CommentTok{\# Calificaciones de riesgo}
\NormalTok{calificaciones }\OperatorTok{=}\NormalTok{ [}\StringTok{\textquotesingle{}AAA\textquotesingle{}}\NormalTok{, }\StringTok{\textquotesingle{}AA\textquotesingle{}}\NormalTok{, }\StringTok{\textquotesingle{}A\textquotesingle{}}\NormalTok{, }\StringTok{\textquotesingle{}BBB\textquotesingle{}}\NormalTok{, }\StringTok{\textquotesingle{}AA\textquotesingle{}}\NormalTok{, }\StringTok{\textquotesingle{}AAA\textquotesingle{}}\NormalTok{, }\StringTok{\textquotesingle{}A\textquotesingle{}}\NormalTok{, }\StringTok{\textquotesingle{}AAA\textquotesingle{}}\NormalTok{]}

\CommentTok{\# Contar ocurrencias}
\NormalTok{num\_aaa }\OperatorTok{=}\NormalTok{ calificaciones.count(}\StringTok{\textquotesingle{}AAA\textquotesingle{}}\NormalTok{)}
\NormalTok{num\_aa }\OperatorTok{=}\NormalTok{ calificaciones.count(}\StringTok{\textquotesingle{}AA\textquotesingle{}}\NormalTok{)}
\NormalTok{num\_a }\OperatorTok{=}\NormalTok{ calificaciones.count(}\StringTok{\textquotesingle{}A\textquotesingle{}}\NormalTok{)}

\BuiltInTok{print}\NormalTok{(}\SpecialStringTok{f"Calificación AAA: }\SpecialCharTok{\{}\NormalTok{num\_aaa}\SpecialCharTok{\}}\SpecialStringTok{ veces"}\NormalTok{)  }\CommentTok{\# 3}
\BuiltInTok{print}\NormalTok{(}\SpecialStringTok{f"Calificación AA: }\SpecialCharTok{\{}\NormalTok{num\_aa}\SpecialCharTok{\}}\SpecialStringTok{ veces"}\NormalTok{)    }\CommentTok{\# 2}
\BuiltInTok{print}\NormalTok{(}\SpecialStringTok{f"Calificación A: }\SpecialCharTok{\{}\NormalTok{num\_a}\SpecialCharTok{\}}\SpecialStringTok{ veces"}\NormalTok{)      }\CommentTok{\# 2}
\end{Highlighting}
\end{Shaded}

\subsection{Método sort()}\label{muxe9todo-sort}

Ordena la lista in-place (modifica la lista original):

\begin{Shaded}
\begin{Highlighting}[]
\CommentTok{\# Tasas de interés desordenadas}
\NormalTok{tasas }\OperatorTok{=}\NormalTok{ [}\FloatTok{0.065}\NormalTok{, }\FloatTok{0.045}\NormalTok{, }\FloatTok{0.075}\NormalTok{, }\FloatTok{0.055}\NormalTok{, }\FloatTok{0.050}\NormalTok{]}
\BuiltInTok{print}\NormalTok{(}\SpecialStringTok{f"Tasas desordenadas: }\SpecialCharTok{\{}\NormalTok{tasas}\SpecialCharTok{\}}\SpecialStringTok{"}\NormalTok{)}

\CommentTok{\# Ordenar ascendentemente}
\NormalTok{tasas.sort()}
\BuiltInTok{print}\NormalTok{(}\SpecialStringTok{f"Tasas ordenadas (asc): }\SpecialCharTok{\{}\NormalTok{tasas}\SpecialCharTok{\}}\SpecialStringTok{"}\NormalTok{)}
\CommentTok{\# [0.045, 0.05, 0.055, 0.065, 0.075]}

\CommentTok{\# Ordenar descendentemente}
\NormalTok{tasas.sort(reverse}\OperatorTok{=}\VariableTok{True}\NormalTok{)}
\BuiltInTok{print}\NormalTok{(}\SpecialStringTok{f"Tasas ordenadas (desc): }\SpecialCharTok{\{}\NormalTok{tasas}\SpecialCharTok{\}}\SpecialStringTok{"}\NormalTok{)}
\CommentTok{\# [0.075, 0.065, 0.055, 0.05, 0.045]}

\CommentTok{\# Ordenar strings}
\NormalTok{paises }\OperatorTok{=}\NormalTok{ [}\StringTok{\textquotesingle{}Perú\textquotesingle{}}\NormalTok{, }\StringTok{\textquotesingle{}Argentina\textquotesingle{}}\NormalTok{, }\StringTok{\textquotesingle{}Chile\textquotesingle{}}\NormalTok{, }\StringTok{\textquotesingle{}Brasil\textquotesingle{}}\NormalTok{, }\StringTok{\textquotesingle{}Colombia\textquotesingle{}}\NormalTok{]}
\NormalTok{paises.sort()}
\BuiltInTok{print}\NormalTok{(}\SpecialStringTok{f"Países ordenados: }\SpecialCharTok{\{}\NormalTok{paises}\SpecialCharTok{\}}\SpecialStringTok{"}\NormalTok{)}
\CommentTok{\# [\textquotesingle{}Argentina\textquotesingle{}, \textquotesingle{}Brasil\textquotesingle{}, \textquotesingle{}Chile\textquotesingle{}, \textquotesingle{}Colombia\textquotesingle{}, \textquotesingle{}Perú\textquotesingle{}]}
\end{Highlighting}
\end{Shaded}

\subsection{Función sorted()}\label{funciuxf3n-sorted}

Retorna una nueva lista ordenada (no modifica la original):

\begin{Shaded}
\begin{Highlighting}[]
\CommentTok{\# PIB de regiones}
\NormalTok{pib\_regiones }\OperatorTok{=}\NormalTok{ [}\DecValTok{100000}\NormalTok{, }\DecValTok{50000}\NormalTok{, }\DecValTok{75000}\NormalTok{, }\DecValTok{30000}\NormalTok{, }\DecValTok{45000}\NormalTok{]}
\BuiltInTok{print}\NormalTok{(}\SpecialStringTok{f"PIB original: }\SpecialCharTok{\{}\NormalTok{pib\_regiones}\SpecialCharTok{\}}\SpecialStringTok{"}\NormalTok{)}

\CommentTok{\# Crear lista ordenada sin modificar original}
\NormalTok{pib\_ordenado }\OperatorTok{=} \BuiltInTok{sorted}\NormalTok{(pib\_regiones)}
\BuiltInTok{print}\NormalTok{(}\SpecialStringTok{f"PIB ordenado: }\SpecialCharTok{\{}\NormalTok{pib\_ordenado}\SpecialCharTok{\}}\SpecialStringTok{"}\NormalTok{)}
\BuiltInTok{print}\NormalTok{(}\SpecialStringTok{f"PIB original (sin cambios): }\SpecialCharTok{\{}\NormalTok{pib\_regiones}\SpecialCharTok{\}}\SpecialStringTok{"}\NormalTok{)}

\CommentTok{\# Ordenar en reversa}
\NormalTok{pib\_descendente }\OperatorTok{=} \BuiltInTok{sorted}\NormalTok{(pib\_regiones, reverse}\OperatorTok{=}\VariableTok{True}\NormalTok{)}
\BuiltInTok{print}\NormalTok{(}\SpecialStringTok{f"PIB descendente: }\SpecialCharTok{\{}\NormalTok{pib\_descendente}\SpecialCharTok{\}}\SpecialStringTok{"}\NormalTok{)}
\end{Highlighting}
\end{Shaded}

\subsection{Método reverse()}\label{muxe9todo-reverse}

Invierte el orden de los elementos:

\begin{Shaded}
\begin{Highlighting}[]
\CommentTok{\# Años en orden cronológico}
\NormalTok{anos }\OperatorTok{=}\NormalTok{ [}\DecValTok{2019}\NormalTok{, }\DecValTok{2020}\NormalTok{, }\DecValTok{2021}\NormalTok{, }\DecValTok{2022}\NormalTok{, }\DecValTok{2023}\NormalTok{]}
\BuiltInTok{print}\NormalTok{(}\SpecialStringTok{f"Años cronológicos: }\SpecialCharTok{\{}\NormalTok{anos}\SpecialCharTok{\}}\SpecialStringTok{"}\NormalTok{)}

\CommentTok{\# Invertir orden}
\NormalTok{anos.reverse()}
\BuiltInTok{print}\NormalTok{(}\SpecialStringTok{f"Años invertidos: }\SpecialCharTok{\{}\NormalTok{anos}\SpecialCharTok{\}}\SpecialStringTok{"}\NormalTok{)}
\CommentTok{\# [2023, 2022, 2021, 2020, 2019]}

\CommentTok{\# Alternativa con slicing}
\NormalTok{anos\_originales }\OperatorTok{=}\NormalTok{ [}\DecValTok{2019}\NormalTok{, }\DecValTok{2020}\NormalTok{, }\DecValTok{2021}\NormalTok{, }\DecValTok{2022}\NormalTok{, }\DecValTok{2023}\NormalTok{]}
\NormalTok{anos\_invertidos }\OperatorTok{=}\NormalTok{ anos\_originales[::}\OperatorTok{{-}}\DecValTok{1}\NormalTok{]}
\BuiltInTok{print}\NormalTok{(}\SpecialStringTok{f"Invertidos con slicing: }\SpecialCharTok{\{}\NormalTok{anos\_invertidos}\SpecialCharTok{\}}\SpecialStringTok{"}\NormalTok{)}
\end{Highlighting}
\end{Shaded}

\section{Diccionarios}\label{diccionarios}

Los diccionarios son colecciones de pares clave-valor, desordenadas
(hasta Python 3.6) u ordenadas por inserción (Python 3.7+), mutables e
indexables por claves. Son ideales para representar datos estructurados
en economía.

\subsection{Creación de diccionarios}\label{creaciuxf3n-de-diccionarios}

\begin{Shaded}
\begin{Highlighting}[]
\CommentTok{\# Diccionario vacío}
\NormalTok{pais\_vacio }\OperatorTok{=}\NormalTok{ \{\}}
\BuiltInTok{print}\NormalTok{(}\BuiltInTok{type}\NormalTok{(pais\_vacio))  }\CommentTok{\# \textless{}class \textquotesingle{}dict\textquotesingle{}\textgreater{}}

\CommentTok{\# Diccionario con datos económicos de un país}
\NormalTok{peru }\OperatorTok{=}\NormalTok{ \{}
    \StringTok{\textquotesingle{}nombre\textquotesingle{}}\NormalTok{: }\StringTok{\textquotesingle{}Perú\textquotesingle{}}\NormalTok{,}
    \StringTok{\textquotesingle{}pib\textquotesingle{}}\NormalTok{: }\DecValTok{242632}\NormalTok{,}
    \StringTok{\textquotesingle{}poblacion\textquotesingle{}}\NormalTok{: }\DecValTok{33715471}\NormalTok{,}
    \StringTok{\textquotesingle{}moneda\textquotesingle{}}\NormalTok{: }\StringTok{\textquotesingle{}Sol\textquotesingle{}}\NormalTok{,}
    \StringTok{\textquotesingle{}codigo\_iso\textquotesingle{}}\NormalTok{: }\StringTok{\textquotesingle{}PER\textquotesingle{}}
\NormalTok{\}}
\BuiltInTok{print}\NormalTok{(peru)}

\CommentTok{\# Diccionario con diferentes tipos de valores}
\NormalTok{indicadores }\OperatorTok{=}\NormalTok{ \{}
    \StringTok{\textquotesingle{}pib\textquotesingle{}}\NormalTok{: }\DecValTok{242632}\NormalTok{,              }\CommentTok{\# int}
    \StringTok{\textquotesingle{}crecimiento\textquotesingle{}}\NormalTok{: }\FloatTok{2.7}\NormalTok{,          }\CommentTok{\# float}
    \StringTok{\textquotesingle{}moneda\textquotesingle{}}\NormalTok{: }\StringTok{\textquotesingle{}PEN\textquotesingle{}}\NormalTok{,             }\CommentTok{\# string}
    \StringTok{\textquotesingle{}emergente\textquotesingle{}}\NormalTok{: }\VariableTok{True}\NormalTok{,           }\CommentTok{\# boolean}
    \StringTok{\textquotesingle{}sectores\textquotesingle{}}\NormalTok{: [}\StringTok{\textquotesingle{}Minería\textquotesingle{}}\NormalTok{, }\StringTok{\textquotesingle{}Agricultura\textquotesingle{}}\NormalTok{],  }\CommentTok{\# lista}
    \StringTok{\textquotesingle{}comercio\textquotesingle{}}\NormalTok{: \{}\StringTok{\textquotesingle{}exp\textquotesingle{}}\NormalTok{: }\DecValTok{63200}\NormalTok{, }\StringTok{\textquotesingle{}imp\textquotesingle{}}\NormalTok{: }\DecValTok{53800}\NormalTok{\}  }\CommentTok{\# diccionario anidado}
\NormalTok{\}}
\BuiltInTok{print}\NormalTok{(indicadores)}
\end{Highlighting}
\end{Shaded}

\subsection{Acceso a elementos}\label{acceso-a-elementos}

\begin{Shaded}
\begin{Highlighting}[]
\CommentTok{\# Datos macroeconómicos}
\NormalTok{macro }\OperatorTok{=}\NormalTok{ \{}
    \StringTok{\textquotesingle{}pib\textquotesingle{}}\NormalTok{: }\DecValTok{242632}\NormalTok{,}
    \StringTok{\textquotesingle{}inflacion\textquotesingle{}}\NormalTok{: }\FloatTok{3.5}\NormalTok{,}
    \StringTok{\textquotesingle{}desempleo\textquotesingle{}}\NormalTok{: }\FloatTok{7.2}\NormalTok{,}
    \StringTok{\textquotesingle{}tipo\_cambio\textquotesingle{}}\NormalTok{: }\FloatTok{3.75}
\NormalTok{\}}

\CommentTok{\# Acceder usando corchetes}
\NormalTok{pib\_actual }\OperatorTok{=}\NormalTok{ macro[}\StringTok{\textquotesingle{}pib\textquotesingle{}}\NormalTok{]}
\BuiltInTok{print}\NormalTok{(}\SpecialStringTok{f"PIB: }\SpecialCharTok{\{}\NormalTok{pib\_actual}\SpecialCharTok{\}}\SpecialStringTok{"}\NormalTok{)}

\NormalTok{inflacion\_actual }\OperatorTok{=}\NormalTok{ macro[}\StringTok{\textquotesingle{}inflacion\textquotesingle{}}\NormalTok{]}
\BuiltInTok{print}\NormalTok{(}\SpecialStringTok{f"Inflación: }\SpecialCharTok{\{}\NormalTok{inflacion\_actual}\SpecialCharTok{\}}\SpecialStringTok{\%"}\NormalTok{)}

\CommentTok{\# Método get() (más seguro)}
\NormalTok{desempleo }\OperatorTok{=}\NormalTok{ macro.get(}\StringTok{\textquotesingle{}desempleo\textquotesingle{}}\NormalTok{)}
\BuiltInTok{print}\NormalTok{(}\SpecialStringTok{f"Desempleo: }\SpecialCharTok{\{}\NormalTok{desempleo}\SpecialCharTok{\}}\SpecialStringTok{\%"}\NormalTok{)}

\CommentTok{\# get() con valor por defecto}
\NormalTok{reservas }\OperatorTok{=}\NormalTok{ macro.get(}\StringTok{\textquotesingle{}reservas\textquotesingle{}}\NormalTok{, }\DecValTok{0}\NormalTok{)  }\CommentTok{\# Retorna 0 si no existe}
\BuiltInTok{print}\NormalTok{(}\SpecialStringTok{f"Reservas: }\SpecialCharTok{\{}\NormalTok{reservas}\SpecialCharTok{\}}\SpecialStringTok{"}\NormalTok{)}

\CommentTok{\# Intentar acceder a clave inexistente con []}
\ControlFlowTok{try}\NormalTok{:}
\NormalTok{    deuda }\OperatorTok{=}\NormalTok{ macro[}\StringTok{\textquotesingle{}deuda\textquotesingle{}}\NormalTok{]}
\ControlFlowTok{except} \PreprocessorTok{KeyError}\NormalTok{:}
    \BuiltInTok{print}\NormalTok{(}\StringTok{"Error: \textquotesingle{}deuda\textquotesingle{} no existe en el diccionario"}\NormalTok{)}
\end{Highlighting}
\end{Shaded}

\subsection{Modificación de
diccionarios}\label{modificaciuxf3n-de-diccionarios}

\begin{Shaded}
\begin{Highlighting}[]
\CommentTok{\# Indicadores económicos}
\NormalTok{indicadores }\OperatorTok{=}\NormalTok{ \{}
    \StringTok{\textquotesingle{}pib\textquotesingle{}}\NormalTok{: }\DecValTok{240000}\NormalTok{,}
    \StringTok{\textquotesingle{}inflacion\textquotesingle{}}\NormalTok{: }\FloatTok{3.2}\NormalTok{,}
    \StringTok{\textquotesingle{}desempleo\textquotesingle{}}\NormalTok{: }\FloatTok{7.5}
\NormalTok{\}}
\BuiltInTok{print}\NormalTok{(}\SpecialStringTok{f"Original: }\SpecialCharTok{\{}\NormalTok{indicadores}\SpecialCharTok{\}}\SpecialStringTok{"}\NormalTok{)}

\CommentTok{\# Modificar valores existentes}
\NormalTok{indicadores[}\StringTok{\textquotesingle{}pib\textquotesingle{}}\NormalTok{] }\OperatorTok{=} \DecValTok{242632}
\NormalTok{indicadores[}\StringTok{\textquotesingle{}inflacion\textquotesingle{}}\NormalTok{] }\OperatorTok{=} \FloatTok{3.5}
\BuiltInTok{print}\NormalTok{(}\SpecialStringTok{f"Actualizado: }\SpecialCharTok{\{}\NormalTok{indicadores}\SpecialCharTok{\}}\SpecialStringTok{"}\NormalTok{)}

\CommentTok{\# Agregar nuevas claves}
\NormalTok{indicadores[}\StringTok{\textquotesingle{}tipo\_cambio\textquotesingle{}}\NormalTok{] }\OperatorTok{=} \FloatTok{3.75}
\NormalTok{indicadores[}\StringTok{\textquotesingle{}reservas\textquotesingle{}}\NormalTok{] }\OperatorTok{=} \DecValTok{74500}
\BuiltInTok{print}\NormalTok{(}\SpecialStringTok{f"Expandido: }\SpecialCharTok{\{}\NormalTok{indicadores}\SpecialCharTok{\}}\SpecialStringTok{"}\NormalTok{)}

\CommentTok{\# Eliminar elementos}
\KeywordTok{del}\NormalTok{ indicadores[}\StringTok{\textquotesingle{}desempleo\textquotesingle{}}\NormalTok{]}
\BuiltInTok{print}\NormalTok{(}\SpecialStringTok{f"Después de eliminar: }\SpecialCharTok{\{}\NormalTok{indicadores}\SpecialCharTok{\}}\SpecialStringTok{"}\NormalTok{)}
\end{Highlighting}
\end{Shaded}

\subsection{Verificación de claves}\label{verificaciuxf3n-de-claves}

\begin{Shaded}
\begin{Highlighting}[]
\CommentTok{\# Datos de un sector económico}
\NormalTok{sector }\OperatorTok{=}\NormalTok{ \{}
    \StringTok{\textquotesingle{}nombre\textquotesingle{}}\NormalTok{: }\StringTok{\textquotesingle{}Minería\textquotesingle{}}\NormalTok{,}
    \StringTok{\textquotesingle{}contribucion\_pib\textquotesingle{}}\NormalTok{: }\FloatTok{10.5}\NormalTok{,}
    \StringTok{\textquotesingle{}empleo\textquotesingle{}}\NormalTok{: }\DecValTok{250000}\NormalTok{,}
    \StringTok{\textquotesingle{}exportaciones\textquotesingle{}}\NormalTok{: }\DecValTok{30000}
\NormalTok{\}}

\CommentTok{\# Verificar si una clave existe}
\BuiltInTok{print}\NormalTok{(}\StringTok{\textquotesingle{}nombre\textquotesingle{}} \KeywordTok{in}\NormalTok{ sector)          }\CommentTok{\# True}
\BuiltInTok{print}\NormalTok{(}\StringTok{\textquotesingle{}importaciones\textquotesingle{}} \KeywordTok{in}\NormalTok{ sector)   }\CommentTok{\# False}

\CommentTok{\# Verificar valores (menos eficiente)}
\BuiltInTok{print}\NormalTok{(}\StringTok{\textquotesingle{}Minería\textquotesingle{}} \KeywordTok{in}\NormalTok{ sector.values())  }\CommentTok{\# True}
\BuiltInTok{print}\NormalTok{(}\FloatTok{10.5} \KeywordTok{in}\NormalTok{ sector.values())       }\CommentTok{\# True}

\CommentTok{\# Uso práctico con condicionales}
\ControlFlowTok{if} \StringTok{\textquotesingle{}contribucion\_pib\textquotesingle{}} \KeywordTok{in}\NormalTok{ sector:}
    \BuiltInTok{print}\NormalTok{(}\SpecialStringTok{f"Contribución al PIB: }\SpecialCharTok{\{}\NormalTok{sector[}\StringTok{\textquotesingle{}contribucion\_pib\textquotesingle{}}\NormalTok{]}\SpecialCharTok{\}}\SpecialStringTok{\%"}\NormalTok{)}
\ControlFlowTok{else}\NormalTok{:}
    \BuiltInTok{print}\NormalTok{(}\StringTok{"No hay datos de contribución al PIB"}\NormalTok{)}
\end{Highlighting}
\end{Shaded}

\section{Operaciones con
diccionarios}\label{operaciones-con-diccionarios}

\subsection{Métodos keys(), values(),
items()}\label{muxe9todos-keys-values-items}

\begin{Shaded}
\begin{Highlighting}[]
\CommentTok{\# Exportaciones por sector (millones USD)}
\NormalTok{exportaciones }\OperatorTok{=}\NormalTok{ \{}
    \StringTok{\textquotesingle{}Minería\textquotesingle{}}\NormalTok{: }\DecValTok{30000}\NormalTok{,}
    \StringTok{\textquotesingle{}Agricultura\textquotesingle{}}\NormalTok{: }\DecValTok{7500}\NormalTok{,}
    \StringTok{\textquotesingle{}Pesca\textquotesingle{}}\NormalTok{: }\DecValTok{2800}\NormalTok{,}
    \StringTok{\textquotesingle{}Manufactura\textquotesingle{}}\NormalTok{: }\DecValTok{12500}
\NormalTok{\}}

\CommentTok{\# Obtener todas las claves}
\NormalTok{sectores }\OperatorTok{=}\NormalTok{ exportaciones.keys()}
\BuiltInTok{print}\NormalTok{(}\SpecialStringTok{f"Sectores: }\SpecialCharTok{\{}\BuiltInTok{list}\NormalTok{(sectores)}\SpecialCharTok{\}}\SpecialStringTok{"}\NormalTok{)}
\CommentTok{\# [\textquotesingle{}Minería\textquotesingle{}, \textquotesingle{}Agricultura\textquotesingle{}, \textquotesingle{}Pesca\textquotesingle{}, \textquotesingle{}Manufactura\textquotesingle{}]}

\CommentTok{\# Obtener todos los valores}
\NormalTok{valores }\OperatorTok{=}\NormalTok{ exportaciones.values()}
\BuiltInTok{print}\NormalTok{(}\SpecialStringTok{f"Valores: }\SpecialCharTok{\{}\BuiltInTok{list}\NormalTok{(valores)}\SpecialCharTok{\}}\SpecialStringTok{"}\NormalTok{)}
\CommentTok{\# [30000, 7500, 2800, 12500]}

\CommentTok{\# Calcular total de exportaciones}
\NormalTok{total\_exportaciones }\OperatorTok{=} \BuiltInTok{sum}\NormalTok{(exportaciones.values())}
\BuiltInTok{print}\NormalTok{(}\SpecialStringTok{f"Total exportaciones: USD }\SpecialCharTok{\{}\NormalTok{total\_exportaciones}\SpecialCharTok{:,\}}\SpecialStringTok{ millones"}\NormalTok{)}

\CommentTok{\# Obtener pares clave{-}valor}
\NormalTok{items }\OperatorTok{=}\NormalTok{ exportaciones.items()}
\BuiltInTok{print}\NormalTok{(}\SpecialStringTok{f"Items: }\SpecialCharTok{\{}\BuiltInTok{list}\NormalTok{(items)}\SpecialCharTok{\}}\SpecialStringTok{"}\NormalTok{)}
\CommentTok{\# [(\textquotesingle{}Minería\textquotesingle{}, 30000), (\textquotesingle{}Agricultura\textquotesingle{}, 7500), ...]}

\CommentTok{\# Iterar sobre items}
\BuiltInTok{print}\NormalTok{(}\StringTok{"}\CharTok{\textbackslash{}n}\StringTok{Detalle de exportaciones:"}\NormalTok{)}
\ControlFlowTok{for}\NormalTok{ sector, valor }\KeywordTok{in}\NormalTok{ exportaciones.items():}
\NormalTok{    porcentaje }\OperatorTok{=}\NormalTok{ (valor }\OperatorTok{/}\NormalTok{ total\_exportaciones) }\OperatorTok{*} \DecValTok{100}
    \BuiltInTok{print}\NormalTok{(}\SpecialStringTok{f"  }\SpecialCharTok{\{}\NormalTok{sector}\SpecialCharTok{\}}\SpecialStringTok{: USD }\SpecialCharTok{\{}\NormalTok{valor}\SpecialCharTok{:,\}}\SpecialStringTok{ M (}\SpecialCharTok{\{}\NormalTok{porcentaje}\SpecialCharTok{:.1f\}}\SpecialStringTok{\%)"}\NormalTok{)}
\end{Highlighting}
\end{Shaded}

\subsection{Método update()}\label{muxe9todo-update}

\begin{Shaded}
\begin{Highlighting}[]
\CommentTok{\# Datos base}
\NormalTok{datos\_2022 }\OperatorTok{=}\NormalTok{ \{}
    \StringTok{\textquotesingle{}pib\textquotesingle{}}\NormalTok{: }\DecValTok{240000}\NormalTok{,}
    \StringTok{\textquotesingle{}inflacion\textquotesingle{}}\NormalTok{: }\FloatTok{3.0}\NormalTok{,}
    \StringTok{\textquotesingle{}desempleo\textquotesingle{}}\NormalTok{: }\FloatTok{7.8}
\NormalTok{\}}
\BuiltInTok{print}\NormalTok{(}\SpecialStringTok{f"Datos 2022: }\SpecialCharTok{\{}\NormalTok{datos\_2022}\SpecialCharTok{\}}\SpecialStringTok{"}\NormalTok{)}

\CommentTok{\# Actualizar con nuevos datos}
\NormalTok{nuevos\_datos }\OperatorTok{=}\NormalTok{ \{}
    \StringTok{\textquotesingle{}pib\textquotesingle{}}\NormalTok{: }\DecValTok{242632}\NormalTok{,}
    \StringTok{\textquotesingle{}tipo\_cambio\textquotesingle{}}\NormalTok{: }\FloatTok{3.75}\NormalTok{,}
    \StringTok{\textquotesingle{}reservas\textquotesingle{}}\NormalTok{: }\DecValTok{74500}
\NormalTok{\}}

\NormalTok{datos\_2022.update(nuevos\_datos)}
\BuiltInTok{print}\NormalTok{(}\SpecialStringTok{f"Datos actualizados: }\SpecialCharTok{\{}\NormalTok{datos\_2022}\SpecialCharTok{\}}\SpecialStringTok{"}\NormalTok{)}
\CommentTok{\# \{\textquotesingle{}pib\textquotesingle{}: 242632, \textquotesingle{}inflacion\textquotesingle{}: 3.0, \textquotesingle{}desempleo\textquotesingle{}: 7.8, }
\CommentTok{\#  \textquotesingle{}tipo\_cambio\textquotesingle{}: 3.75, \textquotesingle{}reservas\textquotesingle{}: 74500\}}
\end{Highlighting}
\end{Shaded}

\subsection{Método pop()}\label{muxe9todo-pop-1}

\begin{Shaded}
\begin{Highlighting}[]
\CommentTok{\# Indicadores trimestrales}
\NormalTok{indicadores }\OperatorTok{=}\NormalTok{ \{}
    \StringTok{\textquotesingle{}Q1\textquotesingle{}}\NormalTok{: }\DecValTok{50000}\NormalTok{,}
    \StringTok{\textquotesingle{}Q2\textquotesingle{}}\NormalTok{: }\DecValTok{52000}\NormalTok{,}
    \StringTok{\textquotesingle{}Q3\textquotesingle{}}\NormalTok{: }\DecValTok{51500}\NormalTok{,}
    \StringTok{\textquotesingle{}Q4\textquotesingle{}}\NormalTok{: }\DecValTok{53000}
\NormalTok{\}}
\BuiltInTok{print}\NormalTok{(}\SpecialStringTok{f"Original: }\SpecialCharTok{\{}\NormalTok{indicadores}\SpecialCharTok{\}}\SpecialStringTok{"}\NormalTok{)}

\CommentTok{\# Eliminar y obtener un valor}
\NormalTok{valor\_q1 }\OperatorTok{=}\NormalTok{ indicadores.pop(}\StringTok{\textquotesingle{}Q1\textquotesingle{}}\NormalTok{)}
\BuiltInTok{print}\NormalTok{(}\SpecialStringTok{f"Valor Q1 eliminado: }\SpecialCharTok{\{}\NormalTok{valor\_q1}\SpecialCharTok{\}}\SpecialStringTok{"}\NormalTok{)}
\BuiltInTok{print}\NormalTok{(}\SpecialStringTok{f"Indicadores restantes: }\SpecialCharTok{\{}\NormalTok{indicadores}\SpecialCharTok{\}}\SpecialStringTok{"}\NormalTok{)}

\CommentTok{\# pop() con valor por defecto}
\NormalTok{valor\_q5 }\OperatorTok{=}\NormalTok{ indicadores.pop(}\StringTok{\textquotesingle{}Q5\textquotesingle{}}\NormalTok{, }\DecValTok{0}\NormalTok{)}
\BuiltInTok{print}\NormalTok{(}\SpecialStringTok{f"Valor Q5: }\SpecialCharTok{\{}\NormalTok{valor\_q5}\SpecialCharTok{\}}\SpecialStringTok{"}\NormalTok{)  }\CommentTok{\# 0 (no existe)}
\end{Highlighting}
\end{Shaded}

\subsection{Método copy()}\label{muxe9todo-copy}

\begin{Shaded}
\begin{Highlighting}[]
\ImportTok{import}\NormalTok{ copy}

\CommentTok{\# Datos originales}
\NormalTok{datos\_base }\OperatorTok{=}\NormalTok{ \{}
    \StringTok{\textquotesingle{}pib\textquotesingle{}}\NormalTok{: }\DecValTok{242632}\NormalTok{,}
    \StringTok{\textquotesingle{}inflacion\textquotesingle{}}\NormalTok{: }\FloatTok{3.5}\NormalTok{,}
    \StringTok{\textquotesingle{}sectores\textquotesingle{}}\NormalTok{: [}\StringTok{\textquotesingle{}Minería\textquotesingle{}}\NormalTok{, }\StringTok{\textquotesingle{}Agricultura\textquotesingle{}}\NormalTok{]}
\NormalTok{\}}

\CommentTok{\# Copia superficial}
\NormalTok{datos\_copia }\OperatorTok{=}\NormalTok{ datos\_base.copy()}
\NormalTok{datos\_copia[}\StringTok{\textquotesingle{}pib\textquotesingle{}}\NormalTok{] }\OperatorTok{=} \DecValTok{250000}
\NormalTok{datos\_copia[}\StringTok{\textquotesingle{}sectores\textquotesingle{}}\NormalTok{].append(}\StringTok{\textquotesingle{}Manufactura\textquotesingle{}}\NormalTok{)}

\BuiltInTok{print}\NormalTok{(}\SpecialStringTok{f"Base: }\SpecialCharTok{\{}\NormalTok{datos\_base}\SpecialCharTok{\}}\SpecialStringTok{"}\NormalTok{)}
\BuiltInTok{print}\NormalTok{(}\SpecialStringTok{f"Copia: }\SpecialCharTok{\{}\NormalTok{datos\_copia}\SpecialCharTok{\}}\SpecialStringTok{"}\NormalTok{)}
\CommentTok{\# Nota: \textquotesingle{}sectores\textquotesingle{} se modifica en ambos}

\CommentTok{\# Copia profunda para diccionarios anidados}
\NormalTok{datos\_profunda }\OperatorTok{=}\NormalTok{ copy.deepcopy(datos\_base)}
\NormalTok{datos\_profunda[}\StringTok{\textquotesingle{}sectores\textquotesingle{}}\NormalTok{].append(}\StringTok{\textquotesingle{}Servicios\textquotesingle{}}\NormalTok{)}

\BuiltInTok{print}\NormalTok{(}\SpecialStringTok{f"}\CharTok{\textbackslash{}n}\SpecialStringTok{Base después de deepcopy: }\SpecialCharTok{\{}\NormalTok{datos\_base}\SpecialCharTok{\}}\SpecialStringTok{"}\NormalTok{)}
\BuiltInTok{print}\NormalTok{(}\SpecialStringTok{f"Copia profunda: }\SpecialCharTok{\{}\NormalTok{datos\_profunda}\SpecialCharTok{\}}\SpecialStringTok{"}\NormalTok{)}
\end{Highlighting}
\end{Shaded}

\subsection{Método clear()}\label{muxe9todo-clear}

\begin{Shaded}
\begin{Highlighting}[]
\CommentTok{\# Datos temporales}
\NormalTok{temp }\OperatorTok{=}\NormalTok{ \{}\StringTok{\textquotesingle{}a\textquotesingle{}}\NormalTok{: }\DecValTok{1}\NormalTok{, }\StringTok{\textquotesingle{}b\textquotesingle{}}\NormalTok{: }\DecValTok{2}\NormalTok{, }\StringTok{\textquotesingle{}c\textquotesingle{}}\NormalTok{: }\DecValTok{3}\NormalTok{\}}
\BuiltInTok{print}\NormalTok{(}\SpecialStringTok{f"Original: }\SpecialCharTok{\{}\NormalTok{temp}\SpecialCharTok{\}}\SpecialStringTok{"}\NormalTok{)}

\NormalTok{temp.clear()}
\BuiltInTok{print}\NormalTok{(}\SpecialStringTok{f"Después de clear: }\SpecialCharTok{\{}\NormalTok{temp}\SpecialCharTok{\}}\SpecialStringTok{"}\NormalTok{)  }\CommentTok{\# \{\}}
\end{Highlighting}
\end{Shaded}

\subsection{Diccionarios anidados}\label{diccionarios-anidados}

\begin{Shaded}
\begin{Highlighting}[]
\CommentTok{\# Datos económicos por país}
\NormalTok{america\_latina }\OperatorTok{=}\NormalTok{ \{}
    \StringTok{\textquotesingle{}Peru\textquotesingle{}}\NormalTok{: \{}
        \StringTok{\textquotesingle{}pib\textquotesingle{}}\NormalTok{: }\DecValTok{242632}\NormalTok{,}
        \StringTok{\textquotesingle{}poblacion\textquotesingle{}}\NormalTok{: }\DecValTok{33715471}\NormalTok{,}
        \StringTok{\textquotesingle{}moneda\textquotesingle{}}\NormalTok{: }\StringTok{\textquotesingle{}PEN\textquotesingle{}}\NormalTok{,}
        \StringTok{\textquotesingle{}capital\textquotesingle{}}\NormalTok{: }\StringTok{\textquotesingle{}Lima\textquotesingle{}}
\NormalTok{    \},}
    \StringTok{\textquotesingle{}Chile\textquotesingle{}}\NormalTok{: \{}
        \StringTok{\textquotesingle{}pib\textquotesingle{}}\NormalTok{: }\DecValTok{301000}\NormalTok{,}
        \StringTok{\textquotesingle{}poblacion\textquotesingle{}}\NormalTok{: }\DecValTok{19493184}\NormalTok{,}
        \StringTok{\textquotesingle{}moneda\textquotesingle{}}\NormalTok{: }\StringTok{\textquotesingle{}CLP\textquotesingle{}}\NormalTok{,}
        \StringTok{\textquotesingle{}capital\textquotesingle{}}\NormalTok{: }\StringTok{\textquotesingle{}Santiago\textquotesingle{}}
\NormalTok{    \},}
    \StringTok{\textquotesingle{}Colombia\textquotesingle{}}\NormalTok{: \{}
        \StringTok{\textquotesingle{}pib\textquotesingle{}}\NormalTok{: }\DecValTok{314000}\NormalTok{,}
        \StringTok{\textquotesingle{}poblacion\textquotesingle{}}\NormalTok{: }\DecValTok{51516562}\NormalTok{,}
        \StringTok{\textquotesingle{}moneda\textquotesingle{}}\NormalTok{: }\StringTok{\textquotesingle{}COP\textquotesingle{}}\NormalTok{,}
        \StringTok{\textquotesingle{}capital\textquotesingle{}}\NormalTok{: }\StringTok{\textquotesingle{}Bogotá\textquotesingle{}}
\NormalTok{    \}}
\NormalTok{\}}

\CommentTok{\# Acceder a datos anidados}
\NormalTok{pib\_peru }\OperatorTok{=}\NormalTok{ america\_latina[}\StringTok{\textquotesingle{}Peru\textquotesingle{}}\NormalTok{][}\StringTok{\textquotesingle{}pib\textquotesingle{}}\NormalTok{]}
\BuiltInTok{print}\NormalTok{(}\SpecialStringTok{f"PIB de Perú: USD }\SpecialCharTok{\{}\NormalTok{pib\_peru}\SpecialCharTok{:,\}}\SpecialStringTok{ millones"}\NormalTok{)}

\NormalTok{capital\_chile }\OperatorTok{=}\NormalTok{ america\_latina[}\StringTok{\textquotesingle{}Chile\textquotesingle{}}\NormalTok{][}\StringTok{\textquotesingle{}capital\textquotesingle{}}\NormalTok{]}
\BuiltInTok{print}\NormalTok{(}\SpecialStringTok{f"Capital de Chile: }\SpecialCharTok{\{}\NormalTok{capital\_chile}\SpecialCharTok{\}}\SpecialStringTok{"}\NormalTok{)}

\CommentTok{\# Calcular PIB per cápita}
\ControlFlowTok{for}\NormalTok{ pais, datos }\KeywordTok{in}\NormalTok{ america\_latina.items():}
\NormalTok{    pib\_pc }\OperatorTok{=}\NormalTok{ (datos[}\StringTok{\textquotesingle{}pib\textquotesingle{}}\NormalTok{] }\OperatorTok{/}\NormalTok{ datos[}\StringTok{\textquotesingle{}poblacion\textquotesingle{}}\NormalTok{]) }\OperatorTok{*} \DecValTok{1000000}
    \BuiltInTok{print}\NormalTok{(}\SpecialStringTok{f"PIB per cápita }\SpecialCharTok{\{}\NormalTok{pais}\SpecialCharTok{\}}\SpecialStringTok{: USD }\SpecialCharTok{\{}\NormalTok{pib\_pc}\SpecialCharTok{:,.2f\}}\SpecialStringTok{"}\NormalTok{)}
\end{Highlighting}
\end{Shaded}

\section{Tuplas}\label{tuplas}

Las tuplas son colecciones ordenadas e inmutables. Son más eficientes en
memoria que las listas y se usan cuando los datos no deben cambiar.

\subsection{Creación de tuplas}\label{creaciuxf3n-de-tuplas}

\begin{Shaded}
\begin{Highlighting}[]
\CommentTok{\# Tupla con paréntesis}
\NormalTok{coordenadas }\OperatorTok{=}\NormalTok{ (}\FloatTok{10.5}\NormalTok{, }\FloatTok{15.3}\NormalTok{)}
\BuiltInTok{print}\NormalTok{(coordenadas)}
\BuiltInTok{print}\NormalTok{(}\BuiltInTok{type}\NormalTok{(coordenadas))  }\CommentTok{\# \textless{}class \textquotesingle{}tuple\textquotesingle{}\textgreater{}}

\CommentTok{\# Tupla sin paréntesis (empaquetado de tupla)}
\NormalTok{punto }\OperatorTok{=} \DecValTok{5}\NormalTok{, }\DecValTok{10}
\BuiltInTok{print}\NormalTok{(punto)}
\BuiltInTok{print}\NormalTok{(}\BuiltInTok{type}\NormalTok{(punto))  }\CommentTok{\# \textless{}class \textquotesingle{}tuple\textquotesingle{}\textgreater{}}

\CommentTok{\# Tupla de un elemento (requiere coma)}
\NormalTok{tupla\_un\_elemento }\OperatorTok{=}\NormalTok{ (}\DecValTok{5}\NormalTok{,)}
\BuiltInTok{print}\NormalTok{(}\BuiltInTok{type}\NormalTok{(tupla\_un\_elemento))  }\CommentTok{\# \textless{}class \textquotesingle{}tuple\textquotesingle{}\textgreater{}}

\NormalTok{no\_es\_tupla }\OperatorTok{=}\NormalTok{ (}\DecValTok{5}\NormalTok{)  }\CommentTok{\# Sin coma, es un int}
\BuiltInTok{print}\NormalTok{(}\BuiltInTok{type}\NormalTok{(no\_es\_tupla))  }\CommentTok{\# \textless{}class \textquotesingle{}int\textquotesingle{}\textgreater{}}

\CommentTok{\# Tupla vacía}
\NormalTok{tupla\_vacia }\OperatorTok{=}\NormalTok{ ()}
\BuiltInTok{print}\NormalTok{(tupla\_vacia)}

\CommentTok{\# Crear tupla desde lista}
\NormalTok{lista\_tasas }\OperatorTok{=}\NormalTok{ [}\FloatTok{0.05}\NormalTok{, }\FloatTok{0.06}\NormalTok{, }\FloatTok{0.055}\NormalTok{, }\FloatTok{0.065}\NormalTok{]}
\NormalTok{tupla\_tasas }\OperatorTok{=} \BuiltInTok{tuple}\NormalTok{(lista\_tasas)}
\BuiltInTok{print}\NormalTok{(tupla\_tasas)}
\end{Highlighting}
\end{Shaded}

\subsection{Características de
tuplas}\label{caracteruxedsticas-de-tuplas}

\begin{Shaded}
\begin{Highlighting}[]
\CommentTok{\# Datos macroeconómicos como tupla}
\NormalTok{macro\_2023 }\OperatorTok{=}\NormalTok{ (}\StringTok{\textquotesingle{}Peru\textquotesingle{}}\NormalTok{, }\DecValTok{242632}\NormalTok{, }\FloatTok{3.5}\NormalTok{, }\FloatTok{7.2}\NormalTok{)}

\CommentTok{\# Indexación (igual que listas)}
\NormalTok{pais }\OperatorTok{=}\NormalTok{ macro\_2023[}\DecValTok{0}\NormalTok{]}
\NormalTok{pib }\OperatorTok{=}\NormalTok{ macro\_2023[}\DecValTok{1}\NormalTok{]}
\BuiltInTok{print}\NormalTok{(}\SpecialStringTok{f"País: }\SpecialCharTok{\{}\NormalTok{pais}\SpecialCharTok{\}}\SpecialStringTok{, PIB: }\SpecialCharTok{\{}\NormalTok{pib}\SpecialCharTok{\}}\SpecialStringTok{"}\NormalTok{)}

\CommentTok{\# Slicing}
\NormalTok{indicadores }\OperatorTok{=}\NormalTok{ macro\_2023[}\DecValTok{2}\NormalTok{:]}
\BuiltInTok{print}\NormalTok{(}\SpecialStringTok{f"Indicadores: }\SpecialCharTok{\{}\NormalTok{indicadores}\SpecialCharTok{\}}\SpecialStringTok{"}\NormalTok{)  }\CommentTok{\# (3.5, 7.2)}

\CommentTok{\# Las tuplas son inmutables}
\ControlFlowTok{try}\NormalTok{:}
\NormalTok{    macro\_2023[}\DecValTok{1}\NormalTok{] }\OperatorTok{=} \DecValTok{250000}
\ControlFlowTok{except} \PreprocessorTok{TypeError}\NormalTok{:}
    \BuiltInTok{print}\NormalTok{(}\StringTok{"Error: Las tuplas no se pueden modificar"}\NormalTok{)}

\CommentTok{\# Pero pueden contener objetos mutables}
\NormalTok{tupla\_con\_lista }\OperatorTok{=}\NormalTok{ (}\DecValTok{1}\NormalTok{, }\DecValTok{2}\NormalTok{, [}\DecValTok{3}\NormalTok{, }\DecValTok{4}\NormalTok{])}
\NormalTok{tupla\_con\_lista[}\DecValTok{2}\NormalTok{].append(}\DecValTok{5}\NormalTok{)  }\CommentTok{\# Modificar la lista interna}
\BuiltInTok{print}\NormalTok{(tupla\_con\_lista)  }\CommentTok{\# (1, 2, [3, 4, 5])}
\end{Highlighting}
\end{Shaded}

\subsection{Ventajas de tuplas}\label{ventajas-de-tuplas}

\begin{Shaded}
\begin{Highlighting}[]
\ImportTok{import}\NormalTok{ sys}

\CommentTok{\# Comparar uso de memoria}
\NormalTok{lista\_grande }\OperatorTok{=}\NormalTok{ [i }\ControlFlowTok{for}\NormalTok{ i }\KeywordTok{in} \BuiltInTok{range}\NormalTok{(}\DecValTok{1000}\NormalTok{)]}
\NormalTok{tupla\_grande }\OperatorTok{=} \BuiltInTok{tuple}\NormalTok{(i }\ControlFlowTok{for}\NormalTok{ i }\KeywordTok{in} \BuiltInTok{range}\NormalTok{(}\DecValTok{1000}\NormalTok{))}

\NormalTok{memoria\_lista }\OperatorTok{=}\NormalTok{ sys.getsizeof(lista\_grande)}
\NormalTok{memoria\_tupla }\OperatorTok{=}\NormalTok{ sys.getsizeof(tupla\_grande)}

\BuiltInTok{print}\NormalTok{(}\SpecialStringTok{f"Memoria lista: }\SpecialCharTok{\{}\NormalTok{memoria\_lista}\SpecialCharTok{\}}\SpecialStringTok{ bytes"}\NormalTok{)}
\BuiltInTok{print}\NormalTok{(}\SpecialStringTok{f"Memoria tupla: }\SpecialCharTok{\{}\NormalTok{memoria\_tupla}\SpecialCharTok{\}}\SpecialStringTok{ bytes"}\NormalTok{)}
\BuiltInTok{print}\NormalTok{(}\SpecialStringTok{f"Tupla usa }\SpecialCharTok{\{}\NormalTok{memoria\_lista }\OperatorTok{{-}}\NormalTok{ memoria\_tupla}\SpecialCharTok{\}}\SpecialStringTok{ bytes menos"}\NormalTok{)}

\CommentTok{\# Tuplas como claves de diccionarios}
\CommentTok{\# (listas no pueden ser claves)}
\NormalTok{coordenadas\_ciudades }\OperatorTok{=}\NormalTok{ \{}
\NormalTok{    (}\StringTok{\textquotesingle{}Lima\textquotesingle{}}\NormalTok{, }\StringTok{\textquotesingle{}Peru\textquotesingle{}}\NormalTok{): (}\FloatTok{12.0464}\NormalTok{, }\FloatTok{77.0428}\NormalTok{),}
\NormalTok{    (}\StringTok{\textquotesingle{}Arequipa\textquotesingle{}}\NormalTok{, }\StringTok{\textquotesingle{}Peru\textquotesingle{}}\NormalTok{): (}\FloatTok{16.4090}\NormalTok{, }\FloatTok{71.5375}\NormalTok{),}
\NormalTok{    (}\StringTok{\textquotesingle{}Cusco\textquotesingle{}}\NormalTok{, }\StringTok{\textquotesingle{}Peru\textquotesingle{}}\NormalTok{): (}\FloatTok{13.5319}\NormalTok{, }\FloatTok{71.9675}\NormalTok{)}
\NormalTok{\}}

\NormalTok{lima\_coords }\OperatorTok{=}\NormalTok{ coordenadas\_ciudades[(}\StringTok{\textquotesingle{}Lima\textquotesingle{}}\NormalTok{, }\StringTok{\textquotesingle{}Peru\textquotesingle{}}\NormalTok{)]}
\BuiltInTok{print}\NormalTok{(}\SpecialStringTok{f"Coordenadas de Lima: }\SpecialCharTok{\{}\NormalTok{lima\_coords}\SpecialCharTok{\}}\SpecialStringTok{"}\NormalTok{)}
\end{Highlighting}
\end{Shaded}

\subsection{Desempaquetado de tuplas}\label{desempaquetado-de-tuplas}

\begin{Shaded}
\begin{Highlighting}[]
\CommentTok{\# Datos de una transacción}
\NormalTok{transaccion }\OperatorTok{=}\NormalTok{ (}\StringTok{\textquotesingle{}Compra\textquotesingle{}}\NormalTok{, }\DecValTok{1000}\NormalTok{, }\FloatTok{3.75}\NormalTok{, }\StringTok{\textquotesingle{}USD\textquotesingle{}}\NormalTok{)}

\CommentTok{\# Desempaquetado simple}
\NormalTok{tipo, cantidad, precio, moneda }\OperatorTok{=}\NormalTok{ transaccion}
\BuiltInTok{print}\NormalTok{(}\SpecialStringTok{f"Tipo: }\SpecialCharTok{\{}\NormalTok{tipo}\SpecialCharTok{\}}\SpecialStringTok{"}\NormalTok{)}
\BuiltInTok{print}\NormalTok{(}\SpecialStringTok{f"Cantidad: }\SpecialCharTok{\{}\NormalTok{cantidad}\SpecialCharTok{\}}\SpecialStringTok{"}\NormalTok{)}
\BuiltInTok{print}\NormalTok{(}\SpecialStringTok{f"Precio: }\SpecialCharTok{\{}\NormalTok{precio}\SpecialCharTok{\}}\SpecialStringTok{"}\NormalTok{)}
\BuiltInTok{print}\NormalTok{(}\SpecialStringTok{f"Moneda: }\SpecialCharTok{\{}\NormalTok{moneda}\SpecialCharTok{\}}\SpecialStringTok{"}\NormalTok{)}

\CommentTok{\# Intercambiar valores sin variable temporal}
\NormalTok{a }\OperatorTok{=} \DecValTok{10}
\NormalTok{b }\OperatorTok{=} \DecValTok{20}
\BuiltInTok{print}\NormalTok{(}\SpecialStringTok{f"Antes: a=}\SpecialCharTok{\{}\NormalTok{a}\SpecialCharTok{\}}\SpecialStringTok{, b=}\SpecialCharTok{\{}\NormalTok{b}\SpecialCharTok{\}}\SpecialStringTok{"}\NormalTok{)}

\NormalTok{a, b }\OperatorTok{=}\NormalTok{ b, a  }\CommentTok{\# Intercambio usando tuplas}
\BuiltInTok{print}\NormalTok{(}\SpecialStringTok{f"Después: a=}\SpecialCharTok{\{}\NormalTok{a}\SpecialCharTok{\}}\SpecialStringTok{, b=}\SpecialCharTok{\{}\NormalTok{b}\SpecialCharTok{\}}\SpecialStringTok{"}\NormalTok{)}

\CommentTok{\# Desempaquetado extendido}
\NormalTok{numeros }\OperatorTok{=}\NormalTok{ (}\DecValTok{1}\NormalTok{, }\DecValTok{2}\NormalTok{, }\DecValTok{3}\NormalTok{, }\DecValTok{4}\NormalTok{, }\DecValTok{5}\NormalTok{)}
\NormalTok{primero, }\OperatorTok{*}\NormalTok{medio, ultimo }\OperatorTok{=}\NormalTok{ numeros}
\BuiltInTok{print}\NormalTok{(}\SpecialStringTok{f"Primero: }\SpecialCharTok{\{}\NormalTok{primero}\SpecialCharTok{\}}\SpecialStringTok{"}\NormalTok{)    }\CommentTok{\# 1}
\BuiltInTok{print}\NormalTok{(}\SpecialStringTok{f"Medio: }\SpecialCharTok{\{}\NormalTok{medio}\SpecialCharTok{\}}\SpecialStringTok{"}\NormalTok{)        }\CommentTok{\# [2, 3, 4]}
\BuiltInTok{print}\NormalTok{(}\SpecialStringTok{f"Último: }\SpecialCharTok{\{}\NormalTok{ultimo}\SpecialCharTok{\}}\SpecialStringTok{"}\NormalTok{)      }\CommentTok{\# 5}
\end{Highlighting}
\end{Shaded}

\subsection{Métodos de tuplas}\label{muxe9todos-de-tuplas}

\begin{Shaded}
\begin{Highlighting}[]
\CommentTok{\# Tuplas tienen pocos métodos (porque son inmutables)}
\NormalTok{calificaciones }\OperatorTok{=}\NormalTok{ (}\StringTok{\textquotesingle{}AAA\textquotesingle{}}\NormalTok{, }\StringTok{\textquotesingle{}AA\textquotesingle{}}\NormalTok{, }\StringTok{\textquotesingle{}A\textquotesingle{}}\NormalTok{, }\StringTok{\textquotesingle{}BBB\textquotesingle{}}\NormalTok{, }\StringTok{\textquotesingle{}AA\textquotesingle{}}\NormalTok{, }\StringTok{\textquotesingle{}AAA\textquotesingle{}}\NormalTok{, }\StringTok{\textquotesingle{}A\textquotesingle{}}\NormalTok{, }\StringTok{\textquotesingle{}AAA\textquotesingle{}}\NormalTok{)}

\CommentTok{\# count() {-} contar ocurrencias}
\NormalTok{num\_aaa }\OperatorTok{=}\NormalTok{ calificaciones.count(}\StringTok{\textquotesingle{}AAA\textquotesingle{}}\NormalTok{)}
\BuiltInTok{print}\NormalTok{(}\SpecialStringTok{f"Número de AAA: }\SpecialCharTok{\{}\NormalTok{num\_aaa}\SpecialCharTok{\}}\SpecialStringTok{"}\NormalTok{)  }\CommentTok{\# 3}

\CommentTok{\# index() {-} encontrar índice}
\NormalTok{indice\_bbb }\OperatorTok{=}\NormalTok{ calificaciones.index(}\StringTok{\textquotesingle{}BBB\textquotesingle{}}\NormalTok{)}
\BuiltInTok{print}\NormalTok{(}\SpecialStringTok{f"Índice de BBB: }\SpecialCharTok{\{}\NormalTok{indice\_bbb}\SpecialCharTok{\}}\SpecialStringTok{"}\NormalTok{)  }\CommentTok{\# 3}

\CommentTok{\# Funciones estándar}
\BuiltInTok{print}\NormalTok{(}\SpecialStringTok{f"Longitud: }\SpecialCharTok{\{}\BuiltInTok{len}\NormalTok{(calificaciones)}\SpecialCharTok{\}}\SpecialStringTok{"}\NormalTok{)  }\CommentTok{\# 8}

\CommentTok{\# Tuplas numéricas}
\NormalTok{precios }\OperatorTok{=}\NormalTok{ (}\DecValTok{100}\NormalTok{, }\DecValTok{105}\NormalTok{, }\DecValTok{98}\NormalTok{, }\DecValTok{110}\NormalTok{, }\DecValTok{95}\NormalTok{)}
\BuiltInTok{print}\NormalTok{(}\SpecialStringTok{f"Máximo: }\SpecialCharTok{\{}\BuiltInTok{max}\NormalTok{(precios)}\SpecialCharTok{\}}\SpecialStringTok{"}\NormalTok{)   }\CommentTok{\# 110}
\BuiltInTok{print}\NormalTok{(}\SpecialStringTok{f"Mínimo: }\SpecialCharTok{\{}\BuiltInTok{min}\NormalTok{(precios)}\SpecialCharTok{\}}\SpecialStringTok{"}\NormalTok{)   }\CommentTok{\# 95}
\BuiltInTok{print}\NormalTok{(}\SpecialStringTok{f"Suma: }\SpecialCharTok{\{}\BuiltInTok{sum}\NormalTok{(precios)}\SpecialCharTok{\}}\SpecialStringTok{"}\NormalTok{)     }\CommentTok{\# 508}
\end{Highlighting}
\end{Shaded}

\section{Asignación múltiple con
diccionarios}\label{asignaciuxf3n-muxfaltiple-con-diccionarios}

El desempaquetado funciona con diccionarios también:

\begin{Shaded}
\begin{Highlighting}[]
\CommentTok{\# Desempaquetar claves}
\NormalTok{datos }\OperatorTok{=}\NormalTok{ \{}\StringTok{\textquotesingle{}nombre\textquotesingle{}}\NormalTok{: }\StringTok{\textquotesingle{}Peru\textquotesingle{}}\NormalTok{, }\StringTok{\textquotesingle{}pib\textquotesingle{}}\NormalTok{: }\DecValTok{242632}\NormalTok{, }\StringTok{\textquotesingle{}moneda\textquotesingle{}}\NormalTok{: }\StringTok{\textquotesingle{}PEN\textquotesingle{}}\NormalTok{\}}

\CommentTok{\# Obtener claves}
\NormalTok{nombre, pib, moneda }\OperatorTok{=}\NormalTok{ datos  }\CommentTok{\# Solo las claves}
\BuiltInTok{print}\NormalTok{(}\SpecialStringTok{f"Claves: }\SpecialCharTok{\{}\NormalTok{nombre}\SpecialCharTok{\}}\SpecialStringTok{, }\SpecialCharTok{\{}\NormalTok{pib}\SpecialCharTok{\}}\SpecialStringTok{, }\SpecialCharTok{\{}\NormalTok{moneda}\SpecialCharTok{\}}\SpecialStringTok{"}\NormalTok{)}
\CommentTok{\# nombre, pib, moneda (strings de las claves)}

\CommentTok{\# Obtener valores}
\NormalTok{nombre\_v, pib\_v, moneda\_v }\OperatorTok{=}\NormalTok{ datos.values()}
\BuiltInTok{print}\NormalTok{(}\SpecialStringTok{f"Valores: }\SpecialCharTok{\{}\NormalTok{nombre\_v}\SpecialCharTok{\}}\SpecialStringTok{, }\SpecialCharTok{\{}\NormalTok{pib\_v}\SpecialCharTok{\}}\SpecialStringTok{, }\SpecialCharTok{\{}\NormalTok{moneda\_v}\SpecialCharTok{\}}\SpecialStringTok{"}\NormalTok{)}
\CommentTok{\# Peru, 242632, PEN}

\CommentTok{\# Desempaquetar items (mejor opción)}
\ControlFlowTok{for}\NormalTok{ clave, valor }\KeywordTok{in}\NormalTok{ datos.items():}
    \BuiltInTok{print}\NormalTok{(}\SpecialStringTok{f"}\SpecialCharTok{\{}\NormalTok{clave}\SpecialCharTok{\}}\SpecialStringTok{: }\SpecialCharTok{\{}\NormalTok{valor}\SpecialCharTok{\}}\SpecialStringTok{"}\NormalTok{)}
\end{Highlighting}
\end{Shaded}

\subsection{Retorno múltiple de
funciones}\label{retorno-muxfaltiple-de-funciones}

\begin{Shaded}
\begin{Highlighting}[]
\KeywordTok{def}\NormalTok{ calcular\_estadisticas(valores):}
    \CommentTok{"""Retorna múltiples estadísticas como tupla"""}
\NormalTok{    minimo }\OperatorTok{=} \BuiltInTok{min}\NormalTok{(valores)}
\NormalTok{    maximo }\OperatorTok{=} \BuiltInTok{max}\NormalTok{(valores)}
\NormalTok{    promedio }\OperatorTok{=} \BuiltInTok{sum}\NormalTok{(valores) }\OperatorTok{/} \BuiltInTok{len}\NormalTok{(valores)}
    \ControlFlowTok{return}\NormalTok{ minimo, maximo, promedio}

\CommentTok{\# Datos de ventas}
\NormalTok{ventas }\OperatorTok{=}\NormalTok{ [}\DecValTok{45}\NormalTok{, }\DecValTok{52}\NormalTok{, }\DecValTok{48}\NormalTok{, }\DecValTok{55}\NormalTok{, }\DecValTok{60}\NormalTok{, }\DecValTok{58}\NormalTok{, }\DecValTok{62}\NormalTok{]}

\CommentTok{\# Desempaquetar retorno}
\NormalTok{min\_venta, max\_venta, prom\_venta }\OperatorTok{=}\NormalTok{ calcular\_estadisticas(ventas)}

\BuiltInTok{print}\NormalTok{(}\SpecialStringTok{f"Venta mínima: }\SpecialCharTok{\{}\NormalTok{min\_venta}\SpecialCharTok{\}}\SpecialStringTok{"}\NormalTok{)}
\BuiltInTok{print}\NormalTok{(}\SpecialStringTok{f"Venta máxima: }\SpecialCharTok{\{}\NormalTok{max\_venta}\SpecialCharTok{\}}\SpecialStringTok{"}\NormalTok{)}
\BuiltInTok{print}\NormalTok{(}\SpecialStringTok{f"Venta promedio: }\SpecialCharTok{\{}\NormalTok{prom\_venta}\SpecialCharTok{:.2f\}}\SpecialStringTok{"}\NormalTok{)}
\end{Highlighting}
\end{Shaded}

\section{Aplicaciones económicas}\label{aplicaciones-econuxf3micas}

\subsection{Ejemplo 1: Análisis de cartera de
inversiones}\label{ejemplo-1-anuxe1lisis-de-cartera-de-inversiones}

\begin{Shaded}
\begin{Highlighting}[]
\CommentTok{\# Cartera de inversiones}
\NormalTok{cartera }\OperatorTok{=}\NormalTok{ \{}
    \StringTok{\textquotesingle{}BBVA\textquotesingle{}}\NormalTok{: \{}\StringTok{\textquotesingle{}cantidad\textquotesingle{}}\NormalTok{: }\DecValTok{100}\NormalTok{, }\StringTok{\textquotesingle{}precio\_compra\textquotesingle{}}\NormalTok{: }\FloatTok{25.50}\NormalTok{, }\StringTok{\textquotesingle{}precio\_actual\textquotesingle{}}\NormalTok{: }\FloatTok{27.80}\NormalTok{\},}
    \StringTok{\textquotesingle{}BCP\textquotesingle{}}\NormalTok{: \{}\StringTok{\textquotesingle{}cantidad\textquotesingle{}}\NormalTok{: }\DecValTok{150}\NormalTok{, }\StringTok{\textquotesingle{}precio\_compra\textquotesingle{}}\NormalTok{: }\FloatTok{180.00}\NormalTok{, }\StringTok{\textquotesingle{}precio\_actual\textquotesingle{}}\NormalTok{: }\FloatTok{195.50}\NormalTok{\},}
    \StringTok{\textquotesingle{}Credicorp\textquotesingle{}}\NormalTok{: \{}\StringTok{\textquotesingle{}cantidad\textquotesingle{}}\NormalTok{: }\DecValTok{50}\NormalTok{, }\StringTok{\textquotesingle{}precio\_compra\textquotesingle{}}\NormalTok{: }\FloatTok{150.00}\NormalTok{, }\StringTok{\textquotesingle{}precio\_actual\textquotesingle{}}\NormalTok{: }\FloatTok{148.00}\NormalTok{\}}
\NormalTok{\}}

\BuiltInTok{print}\NormalTok{(}\StringTok{"ANÁLISIS DE CARTERA DE INVERSIONES"}\NormalTok{)}
\BuiltInTok{print}\NormalTok{(}\StringTok{"="} \OperatorTok{*} \DecValTok{60}\NormalTok{)}

\NormalTok{inversion\_total }\OperatorTok{=} \DecValTok{0}
\NormalTok{valor\_actual\_total }\OperatorTok{=} \DecValTok{0}
\NormalTok{ganancia\_total }\OperatorTok{=} \DecValTok{0}

\ControlFlowTok{for}\NormalTok{ accion, datos }\KeywordTok{in}\NormalTok{ cartera.items():}
\NormalTok{    cantidad }\OperatorTok{=}\NormalTok{ datos[}\StringTok{\textquotesingle{}cantidad\textquotesingle{}}\NormalTok{]}
\NormalTok{    precio\_compra }\OperatorTok{=}\NormalTok{ datos[}\StringTok{\textquotesingle{}precio\_compra\textquotesingle{}}\NormalTok{]}
\NormalTok{    precio\_actual }\OperatorTok{=}\NormalTok{ datos[}\StringTok{\textquotesingle{}precio\_actual\textquotesingle{}}\NormalTok{]}
    
\NormalTok{    inversion }\OperatorTok{=}\NormalTok{ cantidad }\OperatorTok{*}\NormalTok{ precio\_compra}
\NormalTok{    valor\_actual }\OperatorTok{=}\NormalTok{ cantidad }\OperatorTok{*}\NormalTok{ precio\_actual}
\NormalTok{    ganancia }\OperatorTok{=}\NormalTok{ valor\_actual }\OperatorTok{{-}}\NormalTok{ inversion}
\NormalTok{    rentabilidad }\OperatorTok{=}\NormalTok{ (ganancia }\OperatorTok{/}\NormalTok{ inversion) }\OperatorTok{*} \DecValTok{100}
    
\NormalTok{    inversion\_total }\OperatorTok{+=}\NormalTok{ inversion}
\NormalTok{    valor\_actual\_total }\OperatorTok{+=}\NormalTok{ valor\_actual}
\NormalTok{    ganancia\_total }\OperatorTok{+=}\NormalTok{ ganancia}
    
    \BuiltInTok{print}\NormalTok{(}\SpecialStringTok{f"}\CharTok{\textbackslash{}n}\SpecialCharTok{\{}\NormalTok{accion}\SpecialCharTok{\}}\SpecialStringTok{:"}\NormalTok{)}
    \BuiltInTok{print}\NormalTok{(}\SpecialStringTok{f"  Cantidad: }\SpecialCharTok{\{}\NormalTok{cantidad}\SpecialCharTok{\}}\SpecialStringTok{ acciones"}\NormalTok{)}
    \BuiltInTok{print}\NormalTok{(}\SpecialStringTok{f"  Precio compra: S/ }\SpecialCharTok{\{}\NormalTok{precio\_compra}\SpecialCharTok{:.2f\}}\SpecialStringTok{"}\NormalTok{)}
    \BuiltInTok{print}\NormalTok{(}\SpecialStringTok{f"  Precio actual: S/ }\SpecialCharTok{\{}\NormalTok{precio\_actual}\SpecialCharTok{:.2f\}}\SpecialStringTok{"}\NormalTok{)}
    \BuiltInTok{print}\NormalTok{(}\SpecialStringTok{f"  Inversión: S/ }\SpecialCharTok{\{}\NormalTok{inversion}\SpecialCharTok{:,.2f\}}\SpecialStringTok{"}\NormalTok{)}
    \BuiltInTok{print}\NormalTok{(}\SpecialStringTok{f"  Valor actual: S/ }\SpecialCharTok{\{}\NormalTok{valor\_actual}\SpecialCharTok{:,.2f\}}\SpecialStringTok{"}\NormalTok{)}
    \BuiltInTok{print}\NormalTok{(}\SpecialStringTok{f"  Ganancia/Pérdida: S/ }\SpecialCharTok{\{}\NormalTok{ganancia}\SpecialCharTok{:,.2f\}}\SpecialStringTok{"}\NormalTok{)}
    \BuiltInTok{print}\NormalTok{(}\SpecialStringTok{f"  Rentabilidad: }\SpecialCharTok{\{}\NormalTok{rentabilidad}\SpecialCharTok{:+.2f\}}\SpecialStringTok{\%"}\NormalTok{)}

\NormalTok{rentabilidad\_total }\OperatorTok{=}\NormalTok{ (ganancia\_total }\OperatorTok{/}\NormalTok{ inversion\_total) }\OperatorTok{*} \DecValTok{100}

\BuiltInTok{print}\NormalTok{(}\StringTok{"}\CharTok{\textbackslash{}n}\StringTok{"} \OperatorTok{+} \StringTok{"="} \OperatorTok{*} \DecValTok{60}\NormalTok{)}
\BuiltInTok{print}\NormalTok{(}\StringTok{"RESUMEN TOTAL:"}\NormalTok{)}
\BuiltInTok{print}\NormalTok{(}\SpecialStringTok{f"  Inversión total: S/ }\SpecialCharTok{\{}\NormalTok{inversion\_total}\SpecialCharTok{:,.2f\}}\SpecialStringTok{"}\NormalTok{)}
\BuiltInTok{print}\NormalTok{(}\SpecialStringTok{f"  Valor actual total: S/ }\SpecialCharTok{\{}\NormalTok{valor\_actual\_total}\SpecialCharTok{:,.2f\}}\SpecialStringTok{"}\NormalTok{)}
\BuiltInTok{print}\NormalTok{(}\SpecialStringTok{f"  Ganancia/Pérdida total: S/ }\SpecialCharTok{\{}\NormalTok{ganancia\_total}\SpecialCharTok{:,.2f\}}\SpecialStringTok{"}\NormalTok{)}
\BuiltInTok{print}\NormalTok{(}\SpecialStringTok{f"  Rentabilidad total: }\SpecialCharTok{\{}\NormalTok{rentabilidad\_total}\SpecialCharTok{:+.2f\}}\SpecialStringTok{\%"}\NormalTok{)}
\end{Highlighting}
\end{Shaded}

\subsection{Ejemplo 2: Análisis de comercio
exterior}\label{ejemplo-2-anuxe1lisis-de-comercio-exterior}

\begin{Shaded}
\begin{Highlighting}[]
\CommentTok{\# Datos de comercio por país (millones USD)}
\NormalTok{comercio\_exterior }\OperatorTok{=}\NormalTok{ \{}
    \StringTok{\textquotesingle{}China\textquotesingle{}}\NormalTok{: \{}\StringTok{\textquotesingle{}exportaciones\textquotesingle{}}\NormalTok{: }\DecValTok{15000}\NormalTok{, }\StringTok{\textquotesingle{}importaciones\textquotesingle{}}\NormalTok{: }\DecValTok{12000}\NormalTok{\},}
    \StringTok{\textquotesingle{}USA\textquotesingle{}}\NormalTok{: \{}\StringTok{\textquotesingle{}exportaciones\textquotesingle{}}\NormalTok{: }\DecValTok{8500}\NormalTok{, }\StringTok{\textquotesingle{}importaciones\textquotesingle{}}\NormalTok{: }\DecValTok{10500}\NormalTok{\},}
    \StringTok{\textquotesingle{}Brasil\textquotesingle{}}\NormalTok{: \{}\StringTok{\textquotesingle{}exportaciones\textquotesingle{}}\NormalTok{: }\DecValTok{2500}\NormalTok{, }\StringTok{\textquotesingle{}importaciones\textquotesingle{}}\NormalTok{: }\DecValTok{3000}\NormalTok{\},}
    \StringTok{\textquotesingle{}Chile\textquotesingle{}}\NormalTok{: \{}\StringTok{\textquotesingle{}exportaciones\textquotesingle{}}\NormalTok{: }\DecValTok{1800}\NormalTok{, }\StringTok{\textquotesingle{}importaciones\textquotesingle{}}\NormalTok{: }\DecValTok{2200}\NormalTok{\}}
\NormalTok{\}}

\BuiltInTok{print}\NormalTok{(}\StringTok{"ANÁLISIS DE COMERCIO EXTERIOR"}\NormalTok{)}
\BuiltInTok{print}\NormalTok{(}\StringTok{"="} \OperatorTok{*} \DecValTok{60}\NormalTok{)}

\NormalTok{total\_exp }\OperatorTok{=} \DecValTok{0}
\NormalTok{total\_imp }\OperatorTok{=} \DecValTok{0}

\NormalTok{resultados }\OperatorTok{=}\NormalTok{ []}

\ControlFlowTok{for}\NormalTok{ pais, datos }\KeywordTok{in}\NormalTok{ comercio\_exterior.items():}
\NormalTok{    exp }\OperatorTok{=}\NormalTok{ datos[}\StringTok{\textquotesingle{}exportaciones\textquotesingle{}}\NormalTok{]}
\NormalTok{    imp }\OperatorTok{=}\NormalTok{ datos[}\StringTok{\textquotesingle{}importaciones\textquotesingle{}}\NormalTok{]}
\NormalTok{    balanza }\OperatorTok{=}\NormalTok{ exp }\OperatorTok{{-}}\NormalTok{ imp}
\NormalTok{    comercio\_total }\OperatorTok{=}\NormalTok{ exp }\OperatorTok{+}\NormalTok{ imp}
    
\NormalTok{    total\_exp }\OperatorTok{+=}\NormalTok{ exp}
\NormalTok{    total\_imp }\OperatorTok{+=}\NormalTok{ imp}
    
\NormalTok{    resultados.append((pais, exp, imp, balanza, comercio\_total))}

\CommentTok{\# Ordenar por volumen de comercio total}
\NormalTok{resultados.sort(key}\OperatorTok{=}\KeywordTok{lambda}\NormalTok{ x: x[}\DecValTok{4}\NormalTok{], reverse}\OperatorTok{=}\VariableTok{True}\NormalTok{)}

\CommentTok{\# Mostrar resultados}
\ControlFlowTok{for}\NormalTok{ pais, exp, imp, balanza, comercio }\KeywordTok{in}\NormalTok{ resultados:}
    \BuiltInTok{print}\NormalTok{(}\SpecialStringTok{f"}\CharTok{\textbackslash{}n}\SpecialCharTok{\{}\NormalTok{pais}\SpecialCharTok{\}}\SpecialStringTok{:"}\NormalTok{)}
    \BuiltInTok{print}\NormalTok{(}\SpecialStringTok{f"  Exportaciones: USD }\SpecialCharTok{\{}\NormalTok{exp}\SpecialCharTok{:,\}}\SpecialStringTok{ M"}\NormalTok{)}
    \BuiltInTok{print}\NormalTok{(}\SpecialStringTok{f"  Importaciones: USD }\SpecialCharTok{\{}\NormalTok{imp}\SpecialCharTok{:,\}}\SpecialStringTok{ M"}\NormalTok{)}
    \BuiltInTok{print}\NormalTok{(}\SpecialStringTok{f"  Balanza comercial: USD }\SpecialCharTok{\{}\NormalTok{balanza}\SpecialCharTok{:+,\}}\SpecialStringTok{ M"}\NormalTok{)}
    \BuiltInTok{print}\NormalTok{(}\SpecialStringTok{f"  Comercio total: USD }\SpecialCharTok{\{}\NormalTok{comercio}\SpecialCharTok{:,\}}\SpecialStringTok{ M"}\NormalTok{)}
    
    \ControlFlowTok{if}\NormalTok{ balanza }\OperatorTok{\textgreater{}} \DecValTok{0}\NormalTok{:}
        \BuiltInTok{print}\NormalTok{(}\SpecialStringTok{f"  Estado: Superávit"}\NormalTok{)}
    \ControlFlowTok{elif}\NormalTok{ balanza }\OperatorTok{\textless{}} \DecValTok{0}\NormalTok{:}
        \BuiltInTok{print}\NormalTok{(}\SpecialStringTok{f"  Estado: Déficit"}\NormalTok{)}
    \ControlFlowTok{else}\NormalTok{:}
        \BuiltInTok{print}\NormalTok{(}\SpecialStringTok{f"  Estado: Equilibrio"}\NormalTok{)}

\NormalTok{balanza\_total }\OperatorTok{=}\NormalTok{ total\_exp }\OperatorTok{{-}}\NormalTok{ total\_imp}

\BuiltInTok{print}\NormalTok{(}\StringTok{"}\CharTok{\textbackslash{}n}\StringTok{"} \OperatorTok{+} \StringTok{"="} \OperatorTok{*} \DecValTok{60}\NormalTok{)}
\BuiltInTok{print}\NormalTok{(}\StringTok{"TOTALES:"}\NormalTok{)}
\BuiltInTok{print}\NormalTok{(}\SpecialStringTok{f"  Exportaciones totales: USD }\SpecialCharTok{\{}\NormalTok{total\_exp}\SpecialCharTok{:,\}}\SpecialStringTok{ M"}\NormalTok{)}
\BuiltInTok{print}\NormalTok{(}\SpecialStringTok{f"  Importaciones totales: USD }\SpecialCharTok{\{}\NormalTok{total\_imp}\SpecialCharTok{:,\}}\SpecialStringTok{ M"}\NormalTok{)}
\BuiltInTok{print}\NormalTok{(}\SpecialStringTok{f"  Balanza comercial total: USD }\SpecialCharTok{\{}\NormalTok{balanza\_total}\SpecialCharTok{:+,\}}\SpecialStringTok{ M"}\NormalTok{)}
\end{Highlighting}
\end{Shaded}

\subsection{Ejemplo 3: Sistema de calificación de
riesgo}\label{ejemplo-3-sistema-de-calificaciuxf3n-de-riesgo}

\begin{Shaded}
\begin{Highlighting}[]
\CommentTok{\# Empresas y sus indicadores financieros}
\NormalTok{empresas }\OperatorTok{=}\NormalTok{ [}
\NormalTok{    (}\StringTok{\textquotesingle{}TechCorp\textquotesingle{}}\NormalTok{, \{}\StringTok{\textquotesingle{}liquidez\textquotesingle{}}\NormalTok{: }\FloatTok{1.8}\NormalTok{, }\StringTok{\textquotesingle{}endeudamiento\textquotesingle{}}\NormalTok{: }\FloatTok{0.45}\NormalTok{, }\StringTok{\textquotesingle{}rentabilidad\textquotesingle{}}\NormalTok{: }\FloatTok{0.15}\NormalTok{\}),}
\NormalTok{    (}\StringTok{\textquotesingle{}IndusPeru\textquotesingle{}}\NormalTok{, \{}\StringTok{\textquotesingle{}liquidez\textquotesingle{}}\NormalTok{: }\FloatTok{1.2}\NormalTok{, }\StringTok{\textquotesingle{}endeudamiento\textquotesingle{}}\NormalTok{: }\FloatTok{0.65}\NormalTok{, }\StringTok{\textquotesingle{}rentabilidad\textquotesingle{}}\NormalTok{: }\FloatTok{0.08}\NormalTok{\}),}
\NormalTok{    (}\StringTok{\textquotesingle{}AgroSur\textquotesingle{}}\NormalTok{, \{}\StringTok{\textquotesingle{}liquidez\textquotesingle{}}\NormalTok{: }\FloatTok{2.5}\NormalTok{, }\StringTok{\textquotesingle{}endeudamiento\textquotesingle{}}\NormalTok{: }\FloatTok{0.30}\NormalTok{, }\StringTok{\textquotesingle{}rentabilidad\textquotesingle{}}\NormalTok{: }\FloatTok{0.12}\NormalTok{\}),}
\NormalTok{    (}\StringTok{\textquotesingle{}MineroX\textquotesingle{}}\NormalTok{, \{}\StringTok{\textquotesingle{}liquidez\textquotesingle{}}\NormalTok{: }\FloatTok{1.5}\NormalTok{, }\StringTok{\textquotesingle{}endeudamiento\textquotesingle{}}\NormalTok{: }\FloatTok{0.55}\NormalTok{, }\StringTok{\textquotesingle{}rentabilidad\textquotesingle{}}\NormalTok{: }\FloatTok{0.18}\NormalTok{\})}
\NormalTok{]}

\KeywordTok{def}\NormalTok{ calificar\_empresa(indicadores):}
    \CommentTok{"""Califica una empresa según sus indicadores"""}
\NormalTok{    liquidez }\OperatorTok{=}\NormalTok{ indicadores[}\StringTok{\textquotesingle{}liquidez\textquotesingle{}}\NormalTok{]}
\NormalTok{    endeudamiento }\OperatorTok{=}\NormalTok{ indicadores[}\StringTok{\textquotesingle{}endeudamiento\textquotesingle{}}\NormalTok{]}
\NormalTok{    rentabilidad }\OperatorTok{=}\NormalTok{ indicadores[}\StringTok{\textquotesingle{}rentabilidad\textquotesingle{}}\NormalTok{]}
    
\NormalTok{    puntos }\OperatorTok{=} \DecValTok{0}
    
    \CommentTok{\# Evaluar liquidez}
    \ControlFlowTok{if}\NormalTok{ liquidez }\OperatorTok{\textgreater{}=} \FloatTok{2.0}\NormalTok{:}
\NormalTok{        puntos }\OperatorTok{+=} \DecValTok{3}
    \ControlFlowTok{elif}\NormalTok{ liquidez }\OperatorTok{\textgreater{}=} \FloatTok{1.5}\NormalTok{:}
\NormalTok{        puntos }\OperatorTok{+=} \DecValTok{2}
    \ControlFlowTok{elif}\NormalTok{ liquidez }\OperatorTok{\textgreater{}=} \FloatTok{1.0}\NormalTok{:}
\NormalTok{        puntos }\OperatorTok{+=} \DecValTok{1}
    
    \CommentTok{\# Evaluar endeudamiento (menor es mejor)}
    \ControlFlowTok{if}\NormalTok{ endeudamiento }\OperatorTok{\textless{}=} \FloatTok{0.3}\NormalTok{:}
\NormalTok{        puntos }\OperatorTok{+=} \DecValTok{3}
    \ControlFlowTok{elif}\NormalTok{ endeudamiento }\OperatorTok{\textless{}=} \FloatTok{0.5}\NormalTok{:}
\NormalTok{        puntos }\OperatorTok{+=} \DecValTok{2}
    \ControlFlowTok{elif}\NormalTok{ endeudamiento }\OperatorTok{\textless{}=} \FloatTok{0.7}\NormalTok{:}
\NormalTok{        puntos }\OperatorTok{+=} \DecValTok{1}
    
    \CommentTok{\# Evaluar rentabilidad}
    \ControlFlowTok{if}\NormalTok{ rentabilidad }\OperatorTok{\textgreater{}=} \FloatTok{0.15}\NormalTok{:}
\NormalTok{        puntos }\OperatorTok{+=} \DecValTok{3}
    \ControlFlowTok{elif}\NormalTok{ rentabilidad }\OperatorTok{\textgreater{}=} \FloatTok{0.10}\NormalTok{:}
\NormalTok{        puntos }\OperatorTok{+=} \DecValTok{2}
    \ControlFlowTok{elif}\NormalTok{ rentabilidad }\OperatorTok{\textgreater{}=} \FloatTok{0.05}\NormalTok{:}
\NormalTok{        puntos }\OperatorTok{+=} \DecValTok{1}
    
    \CommentTok{\# Determinar calificación}
    \ControlFlowTok{if}\NormalTok{ puntos }\OperatorTok{\textgreater{}=} \DecValTok{8}\NormalTok{:}
        \ControlFlowTok{return} \StringTok{\textquotesingle{}AAA\textquotesingle{}}\NormalTok{, puntos}
    \ControlFlowTok{elif}\NormalTok{ puntos }\OperatorTok{\textgreater{}=} \DecValTok{6}\NormalTok{:}
        \ControlFlowTok{return} \StringTok{\textquotesingle{}AA\textquotesingle{}}\NormalTok{, puntos}
    \ControlFlowTok{elif}\NormalTok{ puntos }\OperatorTok{\textgreater{}=} \DecValTok{4}\NormalTok{:}
        \ControlFlowTok{return} \StringTok{\textquotesingle{}A\textquotesingle{}}\NormalTok{, puntos}
    \ControlFlowTok{elif}\NormalTok{ puntos }\OperatorTok{\textgreater{}=} \DecValTok{2}\NormalTok{:}
        \ControlFlowTok{return} \StringTok{\textquotesingle{}BBB\textquotesingle{}}\NormalTok{, puntos}
    \ControlFlowTok{else}\NormalTok{:}
        \ControlFlowTok{return} \StringTok{\textquotesingle{}BB\textquotesingle{}}\NormalTok{, puntos}

\BuiltInTok{print}\NormalTok{(}\StringTok{"SISTEMA DE CALIFICACIÓN DE RIESGO"}\NormalTok{)}
\BuiltInTok{print}\NormalTok{(}\StringTok{"="} \OperatorTok{*} \DecValTok{70}\NormalTok{)}

\NormalTok{calificaciones }\OperatorTok{=}\NormalTok{ []}

\ControlFlowTok{for}\NormalTok{ empresa, indicadores }\KeywordTok{in}\NormalTok{ empresas:}
\NormalTok{    calificacion, puntos }\OperatorTok{=}\NormalTok{ calificar\_empresa(indicadores)}
\NormalTok{    calificaciones.append((empresa, calificacion, puntos, indicadores))}
    
    \BuiltInTok{print}\NormalTok{(}\SpecialStringTok{f"}\CharTok{\textbackslash{}n}\SpecialCharTok{\{}\NormalTok{empresa}\SpecialCharTok{\}}\SpecialStringTok{:"}\NormalTok{)}
    \BuiltInTok{print}\NormalTok{(}\SpecialStringTok{f"  Liquidez: }\SpecialCharTok{\{}\NormalTok{indicadores[}\StringTok{\textquotesingle{}liquidez\textquotesingle{}}\NormalTok{]}\SpecialCharTok{:.2f\}}\SpecialStringTok{"}\NormalTok{)}
    \BuiltInTok{print}\NormalTok{(}\SpecialStringTok{f"  Endeudamiento: }\SpecialCharTok{\{}\NormalTok{indicadores[}\StringTok{\textquotesingle{}endeudamiento\textquotesingle{}}\NormalTok{]}\SpecialCharTok{:.2\%\}}\SpecialStringTok{"}\NormalTok{)}
    \BuiltInTok{print}\NormalTok{(}\SpecialStringTok{f"  Rentabilidad: }\SpecialCharTok{\{}\NormalTok{indicadores[}\StringTok{\textquotesingle{}rentabilidad\textquotesingle{}}\NormalTok{]}\SpecialCharTok{:.2\%\}}\SpecialStringTok{"}\NormalTok{)}
    \BuiltInTok{print}\NormalTok{(}\SpecialStringTok{f"  Puntos: }\SpecialCharTok{\{}\NormalTok{puntos}\SpecialCharTok{\}}\SpecialStringTok{/9"}\NormalTok{)}
    \BuiltInTok{print}\NormalTok{(}\SpecialStringTok{f"  Calificación: }\SpecialCharTok{\{}\NormalTok{calificacion}\SpecialCharTok{\}}\SpecialStringTok{"}\NormalTok{)}

\CommentTok{\# Ordenar por calificación}
\NormalTok{calificaciones.sort(key}\OperatorTok{=}\KeywordTok{lambda}\NormalTok{ x: x[}\DecValTok{2}\NormalTok{], reverse}\OperatorTok{=}\VariableTok{True}\NormalTok{)}

\BuiltInTok{print}\NormalTok{(}\StringTok{"}\CharTok{\textbackslash{}n}\StringTok{"} \OperatorTok{+} \StringTok{"="} \OperatorTok{*} \DecValTok{70}\NormalTok{)}
\BuiltInTok{print}\NormalTok{(}\StringTok{"RANKING DE EMPRESAS:"}\NormalTok{)}
\ControlFlowTok{for}\NormalTok{ i, (empresa, calif, puntos, \_) }\KeywordTok{in} \BuiltInTok{enumerate}\NormalTok{(calificaciones, }\DecValTok{1}\NormalTok{):}
    \BuiltInTok{print}\NormalTok{(}\SpecialStringTok{f"}\SpecialCharTok{\{}\NormalTok{i}\SpecialCharTok{\}}\SpecialStringTok{. }\SpecialCharTok{\{}\NormalTok{empresa}\SpecialCharTok{\}}\SpecialStringTok{: }\SpecialCharTok{\{}\NormalTok{calif}\SpecialCharTok{\}}\SpecialStringTok{ (}\SpecialCharTok{\{}\NormalTok{puntos}\SpecialCharTok{\}}\SpecialStringTok{ puntos)"}\NormalTok{)}
\end{Highlighting}
\end{Shaded}

\section{Ejercicios prácticos}\label{ejercicios-pruxe1cticos}

\subsection{Ejercicio 1: Administrador de presupuesto
familiar}\label{ejercicio-1-administrador-de-presupuesto-familiar}

\begin{Shaded}
\begin{Highlighting}[]
\CommentTok{\# Sistema de presupuesto mensual}
\NormalTok{presupuesto }\OperatorTok{=}\NormalTok{ \{}
    \StringTok{\textquotesingle{}Alimentos\textquotesingle{}}\NormalTok{: }\DecValTok{1200}\NormalTok{,}
    \StringTok{\textquotesingle{}Transporte\textquotesingle{}}\NormalTok{: }\DecValTok{300}\NormalTok{,}
    \StringTok{\textquotesingle{}Vivienda\textquotesingle{}}\NormalTok{: }\DecValTok{800}\NormalTok{,}
    \StringTok{\textquotesingle{}Servicios\textquotesingle{}}\NormalTok{: }\DecValTok{250}\NormalTok{,}
    \StringTok{\textquotesingle{}Educación\textquotesingle{}}\NormalTok{: }\DecValTok{400}\NormalTok{,}
    \StringTok{\textquotesingle{}Salud\textquotesingle{}}\NormalTok{: }\DecValTok{200}\NormalTok{,}
    \StringTok{\textquotesingle{}Entretenimiento\textquotesingle{}}\NormalTok{: }\DecValTok{150}
\NormalTok{\}}

\NormalTok{gastos\_reales }\OperatorTok{=}\NormalTok{ \{}
    \StringTok{\textquotesingle{}Alimentos\textquotesingle{}}\NormalTok{: }\DecValTok{1250}\NormalTok{,}
    \StringTok{\textquotesingle{}Transporte\textquotesingle{}}\NormalTok{: }\DecValTok{280}\NormalTok{,}
    \StringTok{\textquotesingle{}Vivienda\textquotesingle{}}\NormalTok{: }\DecValTok{800}\NormalTok{,}
    \StringTok{\textquotesingle{}Servicios\textquotesingle{}}\NormalTok{: }\DecValTok{270}\NormalTok{,}
    \StringTok{\textquotesingle{}Educación\textquotesingle{}}\NormalTok{: }\DecValTok{400}\NormalTok{,}
    \StringTok{\textquotesingle{}Salud\textquotesingle{}}\NormalTok{: }\DecValTok{180}\NormalTok{,}
    \StringTok{\textquotesingle{}Entretenimiento\textquotesingle{}}\NormalTok{: }\DecValTok{200}
\NormalTok{\}}

\BuiltInTok{print}\NormalTok{(}\StringTok{"ANÁLISIS DE PRESUPUESTO FAMILIAR"}\NormalTok{)}
\BuiltInTok{print}\NormalTok{(}\StringTok{"="} \OperatorTok{*} \DecValTok{70}\NormalTok{)}

\NormalTok{presupuesto\_total }\OperatorTok{=} \BuiltInTok{sum}\NormalTok{(presupuesto.values())}
\NormalTok{gastos\_totales }\OperatorTok{=} \BuiltInTok{sum}\NormalTok{(gastos\_reales.values())}

\BuiltInTok{print}\NormalTok{(}\SpecialStringTok{f"}\CharTok{\textbackslash{}n}\SpecialStringTok{Presupuesto total: S/ }\SpecialCharTok{\{}\NormalTok{presupuesto\_total}\SpecialCharTok{:,.2f\}}\SpecialStringTok{"}\NormalTok{)}
\BuiltInTok{print}\NormalTok{(}\SpecialStringTok{f"Gastos totales: S/ }\SpecialCharTok{\{}\NormalTok{gastos\_totales}\SpecialCharTok{:,.2f\}}\SpecialStringTok{"}\NormalTok{)}
\BuiltInTok{print}\NormalTok{(}\SpecialStringTok{f"Diferencia: S/ }\SpecialCharTok{\{}\NormalTok{presupuesto\_total }\OperatorTok{{-}}\NormalTok{ gastos\_totales}\SpecialCharTok{:+,.2f\}}\SpecialStringTok{"}\NormalTok{)}

\BuiltInTok{print}\NormalTok{(}\StringTok{"}\CharTok{\textbackslash{}n}\StringTok{DETALLE POR CATEGORÍA:"}\NormalTok{)}
\BuiltInTok{print}\NormalTok{(}\StringTok{"{-}"} \OperatorTok{*} \DecValTok{70}\NormalTok{)}

\NormalTok{categorias\_excedidas }\OperatorTok{=}\NormalTok{ []}
\NormalTok{categorias\_dentro }\OperatorTok{=}\NormalTok{ []}
\NormalTok{ahorros\_por\_categoria }\OperatorTok{=}\NormalTok{ []}

\ControlFlowTok{for}\NormalTok{ categoria }\KeywordTok{in}\NormalTok{ presupuesto.keys():}
\NormalTok{    presu }\OperatorTok{=}\NormalTok{ presupuesto[categoria]}
\NormalTok{    gasto }\OperatorTok{=}\NormalTok{ gastos\_reales[categoria]}
\NormalTok{    diferencia }\OperatorTok{=}\NormalTok{ presu }\OperatorTok{{-}}\NormalTok{ gasto}
\NormalTok{    porcentaje }\OperatorTok{=}\NormalTok{ (gasto }\OperatorTok{/}\NormalTok{ presu) }\OperatorTok{*} \DecValTok{100}
    
    \BuiltInTok{print}\NormalTok{(}\SpecialStringTok{f"}\CharTok{\textbackslash{}n}\SpecialCharTok{\{}\NormalTok{categoria}\SpecialCharTok{\}}\SpecialStringTok{:"}\NormalTok{)}
    \BuiltInTok{print}\NormalTok{(}\SpecialStringTok{f"  Presupuestado: S/ }\SpecialCharTok{\{}\NormalTok{presu}\SpecialCharTok{:,.2f\}}\SpecialStringTok{"}\NormalTok{)}
    \BuiltInTok{print}\NormalTok{(}\SpecialStringTok{f"  Gastado: S/ }\SpecialCharTok{\{}\NormalTok{gasto}\SpecialCharTok{:,.2f\}}\SpecialStringTok{"}\NormalTok{)}
    \BuiltInTok{print}\NormalTok{(}\SpecialStringTok{f"  Diferencia: S/ }\SpecialCharTok{\{}\NormalTok{diferencia}\SpecialCharTok{:+,.2f\}}\SpecialStringTok{"}\NormalTok{)}
    \BuiltInTok{print}\NormalTok{(}\SpecialStringTok{f"  Ejecución: }\SpecialCharTok{\{}\NormalTok{porcentaje}\SpecialCharTok{:.1f\}}\SpecialStringTok{\%"}\NormalTok{)}
    
    \ControlFlowTok{if}\NormalTok{ diferencia }\OperatorTok{\textless{}} \DecValTok{0}\NormalTok{:}
        \BuiltInTok{print}\NormalTok{(}\SpecialStringTok{f"  Estado: ⚠️ EXCEDIDO"}\NormalTok{)}
\NormalTok{        categorias\_excedidas.append((categoria, }\BuiltInTok{abs}\NormalTok{(diferencia)))}
    \ControlFlowTok{else}\NormalTok{:}
        \BuiltInTok{print}\NormalTok{(}\SpecialStringTok{f"  Estado: ✓ Dentro del presupuesto"}\NormalTok{)}
\NormalTok{        categorias\_dentro.append((categoria, diferencia))}
\NormalTok{        ahorros\_por\_categoria.append(diferencia)}

\BuiltInTok{print}\NormalTok{(}\StringTok{"}\CharTok{\textbackslash{}n}\StringTok{"} \OperatorTok{+} \StringTok{"="} \OperatorTok{*} \DecValTok{70}\NormalTok{)}
\BuiltInTok{print}\NormalTok{(}\StringTok{"RESUMEN:"}\NormalTok{)}

\ControlFlowTok{if}\NormalTok{ categorias\_excedidas:}
    \BuiltInTok{print}\NormalTok{(}\StringTok{"}\CharTok{\textbackslash{}n}\StringTok{Categorías excedidas:"}\NormalTok{)}
    \ControlFlowTok{for}\NormalTok{ cat, exceso }\KeywordTok{in}\NormalTok{ categorias\_excedidas:}
        \BuiltInTok{print}\NormalTok{(}\SpecialStringTok{f"  {-} }\SpecialCharTok{\{}\NormalTok{cat}\SpecialCharTok{\}}\SpecialStringTok{: S/ }\SpecialCharTok{\{}\NormalTok{exceso}\SpecialCharTok{:,.2f\}}\SpecialStringTok{ de exceso"}\NormalTok{)}

\ControlFlowTok{if}\NormalTok{ categorias\_dentro:}
    \BuiltInTok{print}\NormalTok{(}\StringTok{"}\CharTok{\textbackslash{}n}\StringTok{Categorías con ahorro:"}\NormalTok{)}
    \ControlFlowTok{for}\NormalTok{ cat, ahorro }\KeywordTok{in}\NormalTok{ categorias\_dentro:}
        \BuiltInTok{print}\NormalTok{(}\SpecialStringTok{f"  {-} }\SpecialCharTok{\{}\NormalTok{cat}\SpecialCharTok{\}}\SpecialStringTok{: S/ }\SpecialCharTok{\{}\NormalTok{ahorro}\SpecialCharTok{:,.2f\}}\SpecialStringTok{ de ahorro"}\NormalTok{)}

\ControlFlowTok{if}\NormalTok{ gastos\_totales }\OperatorTok{\textgreater{}}\NormalTok{ presupuesto\_total:}
    \BuiltInTok{print}\NormalTok{(}\SpecialStringTok{f"}\CharTok{\textbackslash{}n}\SpecialStringTok{⚠️ Déficit total: S/ }\SpecialCharTok{\{}\BuiltInTok{abs}\NormalTok{(presupuesto\_total }\OperatorTok{{-}}\NormalTok{ gastos\_totales)}\SpecialCharTok{:,.2f\}}\SpecialStringTok{"}\NormalTok{)}
\ControlFlowTok{else}\NormalTok{:}
    \BuiltInTok{print}\NormalTok{(}\SpecialStringTok{f"}\CharTok{\textbackslash{}n}\SpecialStringTok{✓ Ahorro total: S/ }\SpecialCharTok{\{}\NormalTok{presupuesto\_total }\OperatorTok{{-}}\NormalTok{ gastos\_totales}\SpecialCharTok{:,.2f\}}\SpecialStringTok{"}\NormalTok{)}
\end{Highlighting}
\end{Shaded}

\subsection{Ejercicio 2: Análisis de series de tiempo
económicas}\label{ejercicio-2-anuxe1lisis-de-series-de-tiempo-econuxf3micas}

\begin{Shaded}
\begin{Highlighting}[]
\CommentTok{\# PIB trimestral de múltiples años}
\NormalTok{pib\_trimestral }\OperatorTok{=}\NormalTok{ \{}
    \DecValTok{2021}\NormalTok{: [}\DecValTok{48000}\NormalTok{, }\DecValTok{49500}\NormalTok{, }\DecValTok{50000}\NormalTok{, }\DecValTok{51000}\NormalTok{],}
    \DecValTok{2022}\NormalTok{: [}\DecValTok{51500}\NormalTok{, }\DecValTok{52500}\NormalTok{, }\DecValTok{53000}\NormalTok{, }\DecValTok{54000}\NormalTok{],}
    \DecValTok{2023}\NormalTok{: [}\DecValTok{54500}\NormalTok{, }\DecValTok{55500}\NormalTok{, }\DecValTok{56000}\NormalTok{, }\DecValTok{57000}\NormalTok{]}
\NormalTok{\}}

\BuiltInTok{print}\NormalTok{(}\StringTok{"ANÁLISIS DE SERIES DE TIEMPO {-} PIB TRIMESTRAL"}\NormalTok{)}
\BuiltInTok{print}\NormalTok{(}\StringTok{"="} \OperatorTok{*} \DecValTok{70}\NormalTok{)}

\CommentTok{\# Calcular PIB anual}
\ControlFlowTok{for}\NormalTok{ ano, trimestres }\KeywordTok{in}\NormalTok{ pib\_trimestral.items():}
\NormalTok{    pib\_anual }\OperatorTok{=} \BuiltInTok{sum}\NormalTok{(trimestres)}
    \BuiltInTok{print}\NormalTok{(}\SpecialStringTok{f"}\CharTok{\textbackslash{}n}\SpecialStringTok{Año }\SpecialCharTok{\{}\NormalTok{ano}\SpecialCharTok{\}}\SpecialStringTok{:"}\NormalTok{)}
    \BuiltInTok{print}\NormalTok{(}\SpecialStringTok{f"  PIB anual: S/ }\SpecialCharTok{\{}\NormalTok{pib\_anual}\SpecialCharTok{:,\}}\SpecialStringTok{ millones"}\NormalTok{)}
    
    \CommentTok{\# PIB por trimestre}
    \ControlFlowTok{for}\NormalTok{ i, pib }\KeywordTok{in} \BuiltInTok{enumerate}\NormalTok{(trimestres, }\DecValTok{1}\NormalTok{):}
\NormalTok{        porcentaje }\OperatorTok{=}\NormalTok{ (pib }\OperatorTok{/}\NormalTok{ pib\_anual) }\OperatorTok{*} \DecValTok{100}
        \BuiltInTok{print}\NormalTok{(}\SpecialStringTok{f"  Q}\SpecialCharTok{\{}\NormalTok{i}\SpecialCharTok{\}}\SpecialStringTok{: S/ }\SpecialCharTok{\{}\NormalTok{pib}\SpecialCharTok{:,\}}\SpecialStringTok{ M (}\SpecialCharTok{\{}\NormalTok{porcentaje}\SpecialCharTok{:.1f\}}\SpecialStringTok{\%)"}\NormalTok{)}
    
    \CommentTok{\# Crecimiento intertrimestral}
    \BuiltInTok{print}\NormalTok{(}\SpecialStringTok{f"}\CharTok{\textbackslash{}n}\SpecialStringTok{  Crecimiento intertrimestral:"}\NormalTok{)}
    \ControlFlowTok{for}\NormalTok{ i }\KeywordTok{in} \BuiltInTok{range}\NormalTok{(}\DecValTok{1}\NormalTok{, }\BuiltInTok{len}\NormalTok{(trimestres)):}
\NormalTok{        crecimiento }\OperatorTok{=}\NormalTok{ ((trimestres[i] }\OperatorTok{{-}}\NormalTok{ trimestres[i}\OperatorTok{{-}}\DecValTok{1}\NormalTok{]) }\OperatorTok{/}\NormalTok{ trimestres[i}\OperatorTok{{-}}\DecValTok{1}\NormalTok{]) }\OperatorTok{*} \DecValTok{100}
        \BuiltInTok{print}\NormalTok{(}\SpecialStringTok{f"    Q}\SpecialCharTok{\{}\NormalTok{i}\SpecialCharTok{\}}\SpecialStringTok{ a Q}\SpecialCharTok{\{}\NormalTok{i}\OperatorTok{+}\DecValTok{1}\SpecialCharTok{\}}\SpecialStringTok{: }\SpecialCharTok{\{}\NormalTok{crecimiento}\SpecialCharTok{:+.2f\}}\SpecialStringTok{\%"}\NormalTok{)}

\CommentTok{\# Análisis interanual}
\BuiltInTok{print}\NormalTok{(}\StringTok{"}\CharTok{\textbackslash{}n}\StringTok{"} \OperatorTok{+} \StringTok{"="} \OperatorTok{*} \DecValTok{70}\NormalTok{)}
\BuiltInTok{print}\NormalTok{(}\StringTok{"CRECIMIENTO INTERANUAL:"}\NormalTok{)}

\NormalTok{anos }\OperatorTok{=} \BuiltInTok{sorted}\NormalTok{(pib\_trimestral.keys())}
\ControlFlowTok{for}\NormalTok{ i }\KeywordTok{in} \BuiltInTok{range}\NormalTok{(}\DecValTok{1}\NormalTok{, }\BuiltInTok{len}\NormalTok{(anos)):}
\NormalTok{    ano\_anterior }\OperatorTok{=}\NormalTok{ anos[i}\OperatorTok{{-}}\DecValTok{1}\NormalTok{]}
\NormalTok{    ano\_actual }\OperatorTok{=}\NormalTok{ anos[i]}
    
\NormalTok{    pib\_ant }\OperatorTok{=} \BuiltInTok{sum}\NormalTok{(pib\_trimestral[ano\_anterior])}
\NormalTok{    pib\_act }\OperatorTok{=} \BuiltInTok{sum}\NormalTok{(pib\_trimestral[ano\_actual])}
    
\NormalTok{    crecimiento }\OperatorTok{=}\NormalTok{ ((pib\_act }\OperatorTok{{-}}\NormalTok{ pib\_ant) }\OperatorTok{/}\NormalTok{ pib\_ant) }\OperatorTok{*} \DecValTok{100}
    
    \BuiltInTok{print}\NormalTok{(}\SpecialStringTok{f"}\SpecialCharTok{\{}\NormalTok{ano\_anterior}\SpecialCharTok{\}}\SpecialStringTok{ a }\SpecialCharTok{\{}\NormalTok{ano\_actual}\SpecialCharTok{\}}\SpecialStringTok{: }\SpecialCharTok{\{}\NormalTok{crecimiento}\SpecialCharTok{:+.2f\}}\SpecialStringTok{\%"}\NormalTok{)}

\CommentTok{\# Identificar mejor y peor trimestre}
\NormalTok{todos\_trimestres }\OperatorTok{=}\NormalTok{ []}
\ControlFlowTok{for}\NormalTok{ ano, trimestres }\KeywordTok{in}\NormalTok{ pib\_trimestral.items():}
    \ControlFlowTok{for}\NormalTok{ i, pib }\KeywordTok{in} \BuiltInTok{enumerate}\NormalTok{(trimestres, }\DecValTok{1}\NormalTok{):}
\NormalTok{        todos\_trimestres.append((ano, i, pib))}

\NormalTok{mejor\_trimestre }\OperatorTok{=} \BuiltInTok{max}\NormalTok{(todos\_trimestres, key}\OperatorTok{=}\KeywordTok{lambda}\NormalTok{ x: x[}\DecValTok{2}\NormalTok{])}
\NormalTok{peor\_trimestre }\OperatorTok{=} \BuiltInTok{min}\NormalTok{(todos\_trimestres, key}\OperatorTok{=}\KeywordTok{lambda}\NormalTok{ x: x[}\DecValTok{2}\NormalTok{])}

\BuiltInTok{print}\NormalTok{(}\StringTok{"}\CharTok{\textbackslash{}n}\StringTok{"} \OperatorTok{+} \StringTok{"="} \OperatorTok{*} \DecValTok{70}\NormalTok{)}
\BuiltInTok{print}\NormalTok{(}\SpecialStringTok{f"Mejor trimestre: Q}\SpecialCharTok{\{}\NormalTok{mejor\_trimestre[}\DecValTok{1}\NormalTok{]}\SpecialCharTok{\}}\SpecialStringTok{ }\SpecialCharTok{\{}\NormalTok{mejor\_trimestre[}\DecValTok{0}\NormalTok{]}\SpecialCharTok{\}}\SpecialStringTok{ {-} S/ }\SpecialCharTok{\{}\NormalTok{mejor\_trimestre[}\DecValTok{2}\NormalTok{]}\SpecialCharTok{:,\}}\SpecialStringTok{ M"}\NormalTok{)}
\BuiltInTok{print}\NormalTok{(}\SpecialStringTok{f"Peor trimestre: Q}\SpecialCharTok{\{}\NormalTok{peor\_trimestre[}\DecValTok{1}\NormalTok{]}\SpecialCharTok{\}}\SpecialStringTok{ }\SpecialCharTok{\{}\NormalTok{peor\_trimestre[}\DecValTok{0}\NormalTok{]}\SpecialCharTok{\}}\SpecialStringTok{ {-} S/ }\SpecialCharTok{\{}\NormalTok{peor\_trimestre[}\DecValTok{2}\NormalTok{]}\SpecialCharTok{:,\}}\SpecialStringTok{ M"}\NormalTok{)}
\end{Highlighting}
\end{Shaded}

\subsection{Ejercicio 3: Sistema de
inventario}\label{ejercicio-3-sistema-de-inventario}

\begin{Shaded}
\begin{Highlighting}[]
\CommentTok{\# Inventario de productos}
\NormalTok{inventario }\OperatorTok{=}\NormalTok{ \{}
    \StringTok{\textquotesingle{}PROD001\textquotesingle{}}\NormalTok{: \{}\StringTok{\textquotesingle{}nombre\textquotesingle{}}\NormalTok{: }\StringTok{\textquotesingle{}Laptop\textquotesingle{}}\NormalTok{, }\StringTok{\textquotesingle{}cantidad\textquotesingle{}}\NormalTok{: }\DecValTok{25}\NormalTok{, }\StringTok{\textquotesingle{}precio\textquotesingle{}}\NormalTok{: }\DecValTok{2500}\NormalTok{, }\StringTok{\textquotesingle{}minimo\textquotesingle{}}\NormalTok{: }\DecValTok{10}\NormalTok{\},}
    \StringTok{\textquotesingle{}PROD002\textquotesingle{}}\NormalTok{: \{}\StringTok{\textquotesingle{}nombre\textquotesingle{}}\NormalTok{: }\StringTok{\textquotesingle{}Mouse\textquotesingle{}}\NormalTok{, }\StringTok{\textquotesingle{}cantidad\textquotesingle{}}\NormalTok{: }\DecValTok{100}\NormalTok{, }\StringTok{\textquotesingle{}precio\textquotesingle{}}\NormalTok{: }\DecValTok{25}\NormalTok{, }\StringTok{\textquotesingle{}minimo\textquotesingle{}}\NormalTok{: }\DecValTok{50}\NormalTok{\},}
    \StringTok{\textquotesingle{}PROD003\textquotesingle{}}\NormalTok{: \{}\StringTok{\textquotesingle{}nombre\textquotesingle{}}\NormalTok{: }\StringTok{\textquotesingle{}Teclado\textquotesingle{}}\NormalTok{, }\StringTok{\textquotesingle{}cantidad\textquotesingle{}}\NormalTok{: }\DecValTok{8}\NormalTok{, }\StringTok{\textquotesingle{}precio\textquotesingle{}}\NormalTok{: }\DecValTok{80}\NormalTok{, }\StringTok{\textquotesingle{}minimo\textquotesingle{}}\NormalTok{: }\DecValTok{15}\NormalTok{\},}
    \StringTok{\textquotesingle{}PROD004\textquotesingle{}}\NormalTok{: \{}\StringTok{\textquotesingle{}nombre\textquotesingle{}}\NormalTok{: }\StringTok{\textquotesingle{}Monitor\textquotesingle{}}\NormalTok{, }\StringTok{\textquotesingle{}cantidad\textquotesingle{}}\NormalTok{: }\DecValTok{30}\NormalTok{, }\StringTok{\textquotesingle{}precio\textquotesingle{}}\NormalTok{: }\DecValTok{800}\NormalTok{, }\StringTok{\textquotesingle{}minimo\textquotesingle{}}\NormalTok{: }\DecValTok{20}\NormalTok{\},}
    \StringTok{\textquotesingle{}PROD005\textquotesingle{}}\NormalTok{: \{}\StringTok{\textquotesingle{}nombre\textquotesingle{}}\NormalTok{: }\StringTok{\textquotesingle{}Impresora\textquotesingle{}}\NormalTok{, }\StringTok{\textquotesingle{}cantidad\textquotesingle{}}\NormalTok{: }\DecValTok{5}\NormalTok{, }\StringTok{\textquotesingle{}precio\textquotesingle{}}\NormalTok{: }\DecValTok{450}\NormalTok{, }\StringTok{\textquotesingle{}minimo\textquotesingle{}}\NormalTok{: }\DecValTok{8}\NormalTok{\}}
\NormalTok{\}}

\BuiltInTok{print}\NormalTok{(}\StringTok{"SISTEMA DE GESTIÓN DE INVENTARIO"}\NormalTok{)}
\BuiltInTok{print}\NormalTok{(}\StringTok{"="} \OperatorTok{*} \DecValTok{70}\NormalTok{)}

\NormalTok{valor\_total\_inventario }\OperatorTok{=} \DecValTok{0}
\NormalTok{productos\_bajo\_stock }\OperatorTok{=}\NormalTok{ []}
\NormalTok{productos\_ok }\OperatorTok{=}\NormalTok{ []}

\ControlFlowTok{for}\NormalTok{ codigo, producto }\KeywordTok{in}\NormalTok{ inventario.items():}
\NormalTok{    nombre }\OperatorTok{=}\NormalTok{ producto[}\StringTok{\textquotesingle{}nombre\textquotesingle{}}\NormalTok{]}
\NormalTok{    cantidad }\OperatorTok{=}\NormalTok{ producto[}\StringTok{\textquotesingle{}cantidad\textquotesingle{}}\NormalTok{]}
\NormalTok{    precio }\OperatorTok{=}\NormalTok{ producto[}\StringTok{\textquotesingle{}precio\textquotesingle{}}\NormalTok{]}
\NormalTok{    minimo }\OperatorTok{=}\NormalTok{ producto[}\StringTok{\textquotesingle{}minimo\textquotesingle{}}\NormalTok{]}
    
\NormalTok{    valor\_producto }\OperatorTok{=}\NormalTok{ cantidad }\OperatorTok{*}\NormalTok{ precio}
\NormalTok{    valor\_total\_inventario }\OperatorTok{+=}\NormalTok{ valor\_producto}
    
\NormalTok{    estado }\OperatorTok{=} \StringTok{"✓ Stock OK"} \ControlFlowTok{if}\NormalTok{ cantidad }\OperatorTok{\textgreater{}=}\NormalTok{ minimo }\ControlFlowTok{else} \StringTok{"⚠️ BAJO STOCK"}
    
    \BuiltInTok{print}\NormalTok{(}\SpecialStringTok{f"}\CharTok{\textbackslash{}n}\SpecialStringTok{Código: }\SpecialCharTok{\{}\NormalTok{codigo}\SpecialCharTok{\}}\SpecialStringTok{"}\NormalTok{)}
    \BuiltInTok{print}\NormalTok{(}\SpecialStringTok{f"  Producto: }\SpecialCharTok{\{}\NormalTok{nombre}\SpecialCharTok{\}}\SpecialStringTok{"}\NormalTok{)}
    \BuiltInTok{print}\NormalTok{(}\SpecialStringTok{f"  Cantidad: }\SpecialCharTok{\{}\NormalTok{cantidad}\SpecialCharTok{\}}\SpecialStringTok{ unidades"}\NormalTok{)}
    \BuiltInTok{print}\NormalTok{(}\SpecialStringTok{f"  Precio unitario: S/ }\SpecialCharTok{\{}\NormalTok{precio}\SpecialCharTok{:.2f\}}\SpecialStringTok{"}\NormalTok{)}
    \BuiltInTok{print}\NormalTok{(}\SpecialStringTok{f"  Valor en inventario: S/ }\SpecialCharTok{\{}\NormalTok{valor\_producto}\SpecialCharTok{:,.2f\}}\SpecialStringTok{"}\NormalTok{)}
    \BuiltInTok{print}\NormalTok{(}\SpecialStringTok{f"  Stock mínimo: }\SpecialCharTok{\{}\NormalTok{minimo}\SpecialCharTok{\}}\SpecialStringTok{ unidades"}\NormalTok{)}
    \BuiltInTok{print}\NormalTok{(}\SpecialStringTok{f"  Estado: }\SpecialCharTok{\{}\NormalTok{estado}\SpecialCharTok{\}}\SpecialStringTok{"}\NormalTok{)}
    
    \ControlFlowTok{if}\NormalTok{ cantidad }\OperatorTok{\textless{}}\NormalTok{ minimo:}
\NormalTok{        faltante }\OperatorTok{=}\NormalTok{ minimo }\OperatorTok{{-}}\NormalTok{ cantidad}
\NormalTok{        productos\_bajo\_stock.append((nombre, cantidad, minimo, faltante))}
    \ControlFlowTok{else}\NormalTok{:}
\NormalTok{        productos\_ok.append((nombre, cantidad))}

\BuiltInTok{print}\NormalTok{(}\StringTok{"}\CharTok{\textbackslash{}n}\StringTok{"} \OperatorTok{+} \StringTok{"="} \OperatorTok{*} \DecValTok{70}\NormalTok{)}
\BuiltInTok{print}\NormalTok{(}\StringTok{"RESUMEN DE INVENTARIO:"}\NormalTok{)}
\BuiltInTok{print}\NormalTok{(}\SpecialStringTok{f"Valor total del inventario: S/ }\SpecialCharTok{\{}\NormalTok{valor\_total\_inventario}\SpecialCharTok{:,.2f\}}\SpecialStringTok{"}\NormalTok{)}
\BuiltInTok{print}\NormalTok{(}\SpecialStringTok{f"Productos en stock adecuado: }\SpecialCharTok{\{}\BuiltInTok{len}\NormalTok{(productos\_ok)}\SpecialCharTok{\}}\SpecialStringTok{"}\NormalTok{)}
\BuiltInTok{print}\NormalTok{(}\SpecialStringTok{f"Productos bajo stock mínimo: }\SpecialCharTok{\{}\BuiltInTok{len}\NormalTok{(productos\_bajo\_stock)}\SpecialCharTok{\}}\SpecialStringTok{"}\NormalTok{)}

\ControlFlowTok{if}\NormalTok{ productos\_bajo\_stock:}
    \BuiltInTok{print}\NormalTok{(}\StringTok{"}\CharTok{\textbackslash{}n}\StringTok{⚠️ ALERTA {-} Productos que requieren reposición:"}\NormalTok{)}
    \ControlFlowTok{for}\NormalTok{ nombre, actual, minimo, faltante }\KeywordTok{in}\NormalTok{ productos\_bajo\_stock:}
        \BuiltInTok{print}\NormalTok{(}\SpecialStringTok{f"  {-} }\SpecialCharTok{\{}\NormalTok{nombre}\SpecialCharTok{\}}\SpecialStringTok{: }\SpecialCharTok{\{}\NormalTok{actual}\SpecialCharTok{\}}\SpecialStringTok{ unidades (faltante: }\SpecialCharTok{\{}\NormalTok{faltante}\SpecialCharTok{\}}\SpecialStringTok{ unidades)"}\NormalTok{)}

\BuiltInTok{print}\NormalTok{(}\StringTok{"}\CharTok{\textbackslash{}n}\StringTok{"} \OperatorTok{+} \StringTok{"="} \OperatorTok{*} \DecValTok{70}\NormalTok{)}
\BuiltInTok{print}\NormalTok{(}\StringTok{"PRODUCTOS POR VALOR:"}\NormalTok{)}
\NormalTok{productos\_valor }\OperatorTok{=}\NormalTok{ [(codigo, prod[}\StringTok{\textquotesingle{}nombre\textquotesingle{}}\NormalTok{], prod[}\StringTok{\textquotesingle{}cantidad\textquotesingle{}}\NormalTok{] }\OperatorTok{*}\NormalTok{ prod[}\StringTok{\textquotesingle{}precio\textquotesingle{}}\NormalTok{]) }
                   \ControlFlowTok{for}\NormalTok{ codigo, prod }\KeywordTok{in}\NormalTok{ inventario.items()]}
\NormalTok{productos\_valor.sort(key}\OperatorTok{=}\KeywordTok{lambda}\NormalTok{ x: x[}\DecValTok{2}\NormalTok{], reverse}\OperatorTok{=}\VariableTok{True}\NormalTok{)}

\ControlFlowTok{for}\NormalTok{ i, (codigo, nombre, valor) }\KeywordTok{in} \BuiltInTok{enumerate}\NormalTok{(productos\_valor, }\DecValTok{1}\NormalTok{):}
\NormalTok{    porcentaje }\OperatorTok{=}\NormalTok{ (valor }\OperatorTok{/}\NormalTok{ valor\_total\_inventario) }\OperatorTok{*} \DecValTok{100}
    \BuiltInTok{print}\NormalTok{(}\SpecialStringTok{f"}\SpecialCharTok{\{}\NormalTok{i}\SpecialCharTok{\}}\SpecialStringTok{. }\SpecialCharTok{\{}\NormalTok{nombre}\SpecialCharTok{\}}\SpecialStringTok{ (}\SpecialCharTok{\{}\NormalTok{codigo}\SpecialCharTok{\}}\SpecialStringTok{): S/ }\SpecialCharTok{\{}\NormalTok{valor}\SpecialCharTok{:,.2f\}}\SpecialStringTok{ (}\SpecialCharTok{\{}\NormalTok{porcentaje}\SpecialCharTok{:.1f\}}\SpecialStringTok{\%)"}\NormalTok{)}
\end{Highlighting}
\end{Shaded}

\section{Conclusión}\label{conclusiuxf3n}

En esta guía hemos explorado las estructuras de datos fundamentales de
Python:

\textbf{Listas}

Colecciones ordenadas y mutables. Ideales para secuencias de datos que
pueden cambiar. Métodos importantes: append(), extend(), insert(),
remove(), pop(), sort().

\textbf{Diccionarios}

Colecciones de pares clave-valor. Perfectos para datos estructurados y
búsquedas rápidas. Métodos importantes: keys(), values(), items(),
get(), update(), pop().

\textbf{Tuplas}

Colecciones ordenadas e inmutables. Más eficientes en memoria que las
listas. Útiles para datos que no deben cambiar y como claves de
diccionarios.

\textbf{Conceptos importantes}

La diferencia entre asignación por referencia y por valor es crucial.
Usar copy.copy() para copias superficiales y copy.deepcopy() para copias
profundas de estructuras anidadas.

\section{Próximos pasos}\label{pruxf3ximos-pasos}

En la siguiente guía (Guía 4: Ejecución Condicional) exploraremos:

\begin{itemize}
\tightlist
\item
  Expresiones booleanas
\item
  Operadores lógicos
\item
  Ejecuciones condicionales (if, elif, else)
\item
  Condicionales anidados
\item
  Manejo de excepciones con try-except
\item
  Evaluación de cortocircuito
\end{itemize}

Estos conceptos nos permitirán crear programas que tomen decisiones
basadas en condiciones, esenciales para análisis económico automatizado.

\section{Recursos adicionales}\label{recursos-adicionales}

Para profundizar en estructuras de datos:

\begin{itemize}
\tightlist
\item
  Python Data Structures: docs.python.org/tutorial/datastructures.html
\item
  Práctica con ejercicios económicos reales
\item
  Exploración de bibliotecas como Pandas para análisis avanzado
\end{itemize}

La práctica constante con datos económicos reales consolidará estos
conceptos. Se recomienda trabajar con datos de fuentes como BCRP, INEI,
o Banco Mundial para aplicar estas estructuras en contextos realistas.

\section{Publicaciones Similares}\label{publicaciones-similares}

Si te interesó este artículo, te recomendamos que explores otros blogs y
recursos relacionados que pueden ampliar tus conocimientos. Aquí te dejo
algunas sugerencias:

\begin{enumerate}
\def\labelenumi{\arabic{enumi}.}
\tightlist
\item
  \href{https://numerus-scriptum.netlify.app/python/2020-06-19-instalacion-de-anaconda/index.pdf}{\faIcon{file-pdf}}
  \href{https://numerus-scriptum.netlify.app/python/2020-06-19-instalacion-de-anaconda}{Instalacion
  De Anaconda}
\item
  \href{https://numerus-scriptum.netlify.app/python/2020-06-20-configurar-entorno-virtual-python-anaconda/index.pdf}{\faIcon{file-pdf}}
  \href{https://numerus-scriptum.netlify.app/python/2020-06-20-configurar-entorno-virtual-python-anaconda}{Configurar
  Entorno Virtual Python Anaconda}
\item
  \href{https://numerus-scriptum.netlify.app/python/2021-04-17-01-introducion-a-la-programacion-con-python/index.pdf}{\faIcon{file-pdf}}
  \href{https://numerus-scriptum.netlify.app/python/2021-04-17-01-introducion-a-la-programacion-con-python}{01
  Introducion A La Programacion Con Python}
\item
  \href{https://numerus-scriptum.netlify.app/python/2021-05-31-02-variables-expresiones-y-statements-con-python/index.pdf}{\faIcon{file-pdf}}
  \href{https://numerus-scriptum.netlify.app/python/2021-05-31-02-variables-expresiones-y-statements-con-python}{02
  Variables Expresiones Y Statements Con Python}
\item
  \href{https://numerus-scriptum.netlify.app/python/2021-06-07-03-objetos-de-python/index.pdf}{\faIcon{file-pdf}}
  \href{https://numerus-scriptum.netlify.app/python/2021-06-07-03-objetos-de-python}{03
  Objetos De Python}
\item
  \href{https://numerus-scriptum.netlify.app/python/2021-06-14-04-ejecucion-condicional-con-python/index.pdf}{\faIcon{file-pdf}}
  \href{https://numerus-scriptum.netlify.app/python/2021-06-14-04-ejecucion-condicional-con-python}{04
  Ejecucion Condicional Con Python}
\item
  \href{https://numerus-scriptum.netlify.app/python/2021-06-21-05-iteraciones-con-python/index.pdf}{\faIcon{file-pdf}}
  \href{https://numerus-scriptum.netlify.app/python/2021-06-21-05-iteraciones-con-python}{05
  Iteraciones Con Python}
\item
  \href{https://numerus-scriptum.netlify.app/python/2021-08-16-06-funciones-con-python/index.pdf}{\faIcon{file-pdf}}
  \href{https://numerus-scriptum.netlify.app/python/2021-08-16-06-funciones-con-python}{06
  Funciones Con Python}
\item
  \href{https://numerus-scriptum.netlify.app/python/2021-08-23-07-dataframes-con-python/index.pdf}{\faIcon{file-pdf}}
  \href{https://numerus-scriptum.netlify.app/python/2021-08-23-07-dataframes-con-python}{07
  Dataframes Con Python}
\item
  \href{https://numerus-scriptum.netlify.app/python/2021-11-29-08-prediccion-y-metrica-de-performance-con-python/index.pdf}{\faIcon{file-pdf}}
  \href{https://numerus-scriptum.netlify.app/python/2021-11-29-08-prediccion-y-metrica-de-performance-con-python}{08
  Prediccion Y Metrica De Performance Con Python}
\item
  \href{https://numerus-scriptum.netlify.app/python/2021-12-06-09-metodos-de-machine-learning-para-clasificacion-con-python/index.pdf}{\faIcon{file-pdf}}
  \href{https://numerus-scriptum.netlify.app/python/2021-12-06-09-metodos-de-machine-learning-para-clasificacion-con-python}{09
  Metodos De Machine Learning Para Clasificacion Con Python}
\item
  \href{https://numerus-scriptum.netlify.app/python/2021-12-13-10-metodos-de-machine-learning-para-regresion-con-python/index.pdf}{\faIcon{file-pdf}}
  \href{https://numerus-scriptum.netlify.app/python/2021-12-13-10-metodos-de-machine-learning-para-regresion-con-python}{10
  Metodos De Machine Learning Para Regresion Con Python}
\item
  \href{https://numerus-scriptum.netlify.app/python/2022-10-31-11-validacion-cruzada-y-composicion-del-modelo-con-python/index.pdf}{\faIcon{file-pdf}}
  \href{https://numerus-scriptum.netlify.app/python/2022-10-31-11-validacion-cruzada-y-composicion-del-modelo-con-python}{11
  Validacion Cruzada Y Composicion Del Modelo Con Python}
\item
  \href{https://numerus-scriptum.netlify.app/python/2025-05-10-visualizacion-de-datos-con-python/index.pdf}{\faIcon{file-pdf}}
  \href{https://numerus-scriptum.netlify.app/python/2025-05-10-visualizacion-de-datos-con-python}{Visualizacion
  De Datos Con Python}
\end{enumerate}

Esperamos que encuentres estas publicaciones igualmente interesantes y
útiles. ¡Disfruta de la lectura!






\end{document}
